\documentclass[11pt,a4paper]{article}

\usepackage{amsmath,amssymb,amsthm}
\usepackage{mathtools}
\usepackage{hyperref}
\hypersetup{
  pdftitle={Toward Three Spatial Directions from a Supplied Spinor Carrier: Requirements for an
    Equivariant Bridge},
  pdfauthor={Jérôme Beau},
  pdfsubject={Cosmochrony spectral geometry programme, paper Q7, version 2.0},
}
\usepackage{geometry}
\usepackage{doi}
\geometry{margin=2.5cm}

\newtheorem{theorem}{Theorem}[section]
\newtheorem{lemma}[theorem]{Lemma}
\newtheorem{proposition}[theorem]{Proposition}
\newtheorem{corollary}[theorem]{Corollary}
\newtheorem{definition}[theorem]{Definition}
\newtheorem{remark}[theorem]{Remark}
\newtheorem{conjecture}[theorem]{Conjecture}
\newtheorem{criterion}[theorem]{Criterion}

\newcommand{\Heis}{\mathrm{Heis}_3}
\newcommand{\Heff}{H_{\mathrm{eff}}}
\newcommand{\ImH}{\mathrm{Im}\,\mathbb{H}}
\newcommand{\su}{\mathfrak{su}}
\newcommand{\heis}{\mathfrak{heis}_3}
\newcommand{\Vr}{V_\rho}
\newcommand{\Dhom}{D_{\mathrm{hom}}}
\newcommand{\Sym}{\mathrm{Sym}}

\title{Toward Three Spatial Directions from a Supplied Spinor Carrier:\\
Requirements for an Equivariant Bridge\\[4pt]
\large Q7 of the Cosmochrony Spectral Geometry Programme}
\author{J\'er\^ome Beau\\[4pt]
\small Independent Researcher\\
\small \texttt{jerome.beau@cosmochrony.org}}
\date{2026}

\begin{document}
  \maketitle

  \begin{abstract}
    Two sub-programmes of the Cosmochrony framework produce the integer three as a structural
    output, and this paper states what an equivariant bridge between them would require.
    On the geometric side, paper Q5b derives an effective Lorentzian co-metric on
    $\mathbb{R}_\tau \times \Heis(\mathbb{R})$ whose spatial sector has three directions, the two
    horizontal generators $\tilde X, \tilde Y$ and the central element
    $\tilde Z = [\tilde X, \tilde Y]$; its extraction is conditional on the unestablished spatial
    limit hypothesis [H-L], its central slot is degenerate at principal-symbol order, and the
    Lorentzian signature carries a further hypothesis.
    On the representation side, paper O23 proves that the neutral sector of a \emph{supplied}
    spinor carrier $V_\rho \cong \mathbb{C}^2$ is $\su(2) \simeq \ImH$, of real dimension three,
    and paper O28 works with an admissible projection space $\Heff = \mathbb{C}^3$, supplied as a
    parameter consistent with that selection rule, whose per-pair covariance has rank three as a
    finite-data value of its protocol and not as an invariant.
    The identification of that measured space with a spin-$1$ module of $\su(2)$ is supplied, not
    derived: O29 identifies no such structure, and $d_\rho = 2$ is O26's minimality selection.
    Three requirements are proved here.
    First, the bridge cannot be a Lie algebra isomorphism: $\su(2)$ is semisimple, $\heis$ is
    nilpotent.
    Second, under the identification hypothesis, $\Sym^2(\Vr)$ is, up to isomorphism, the unique
    irreducible three-dimensional complex module of $\su(2)$, so any \emph{non-zero} equivariant
    map to the complexified adjoint module is an isomorphism, unique up to one non-zero complex
    scalar; reality and positivity of that scalar require separate hypotheses on the forms
    compared.
    Third, the $\su(2)$ Casimir acts on $\Sym^2(\Vr)$ as $2 \cdot \mathrm{Id}$, and a \emph{non-zero}
    equivariant bridge pulls the Hermitian spatial form back to a positive-definite form of rank
    three whenever the two coefficients are positive; invariance under the Cartan generator $J_3$ separates the
    spin-weight $\pm 1$
    sector from the spin-weight $0$ sector and
    forces isotropy \emph{within the horizontal plane}, but it does not by itself force
    $A_H = A_z$.
    Version~2.0 re-types the numerical section.
    The quantities previously reported as symbol coefficients are compressions of the discrete
    Weil Laplacian $L_{\mathrm{Weil}} = 4I - W_a - W_a^\dagger - W_b - W_b^\dagger$ onto the three
    stored basis rows, which are pure Fourier modes: the central value is the flat-mode Rayleigh
    quotient $4 - 2 - 0 = 2$, exactly and for every tested pair, and the horizontal value is
    $4 - 2\cos(2\pi k/q)$ for the mode index $k$ selected by the pipeline.
    Consequently the identification of the central value with the Casimir eigenvalue and the
    fitted convergence of the isotropy gap are no longer asserted: for these modes
    $A_H - 2 = 2(1 - \cos(2\pi k/q))$, so vanishing of the gap is equivalent to
    $\mathrm{dist}(k/q, \mathbb{Z}) \to 0$, a property of the mode-selection rule that remains to
    be established.
    The bridge is therefore a missing identification, not a proved obstruction.
    The only obstruction located in the corpus is conditional: under Q5a's depth and weight
    hypotheses, no common scalar normalisation of the published admissibility form yields a toric
    differential limit.
  \end{abstract}

  \noindent\textbf{Keywords:} Heisenberg group; Weil representation; spectral geometry;
  sub-Riemannian Laplacian; Carnot geometry; spin representations; $\mathfrak{su}(2)$ Casimir;
  metaplectic symmetry; Rayleigh quotient; Cosmochrony.

  \clearpage
  \tableofcontents

  \clearpage
  \section{Introduction and Statement of the Problem}

    \subsection{Context}

      The Cosmochrony programme derives physical observables from the spectral properties of
      admissible non-injective projections on the Weil representation of $\Heis(\mathbb{Z}/q\mathbb{Z})$.
      Two sub-programmes have now independently produced the integer~$3$ as a structural output.

      \paragraph{The geometric ``3'' (Q-series).}
        Paper Q5b~\cite{Beau2026q5b} establishes that the BFS shell stratification of
        $G_q = \Heis(\mathbb{Z}/q\mathbb{Z})$ converges, in the pre-saturation regime, to the
        Carnot--Carath\'eodory foliation of $\Heis(\mathbb{R})$.
        The homogeneous dimension $\Dhom = 4$ (Bass--Guivarch) gives the limiting geometry the
        spectral dimension of a four-dimensional space.
        Under the hypotheses of~\cite{Beau2026q5b}, including its unestablished [H-L] and, for the
        signature, a further hyperbolicity hypothesis, the effective co-metric is read with
        signature $(-,+,+,+)$: one temporal direction~$\tau$ and a three-dimensional spatial
        sector.
        Within the sub-Riemannian structure of $\Heis(\mathbb{R})$, these three spatial directions
        arise as: the two horizontal generators $\tilde X, \tilde Y$ (entering the sub-Riemannian
        Laplacian $\Delta_H = \tilde X^2 + \tilde Y^2$) and the central element
        $\tilde Z = [\tilde X, \tilde Y]$, which enters the effective metric via sub-principal
        corrections (open problem Q5b-O2 of~\cite{Beau2026q5b}).

      \paragraph{The admissible ``3'' (O-series).}
        Paper O23~\cite{Beau2026a27} proves that if the admissible neutral sector is carried by an
        irreducible $\mathrm{SU}(2)$-valued two-dimensional representation
        $V_\rho \cong \mathbb{C}^2$ (a supplied input), then its traceless anti-Hermitian sector is
        $\su(2) \simeq \ImH \simeq \mathbb{R}^3$ (Theorem~3.1).
        The carrier selection and the identification of $\Sigma_c$ with that dimension are open
        bridges; $\Sigma_c(n_3) = 3$ enters as a supplied selection rule.
        Paper O28~\cite{Beau2026a32} sets $\Heff = \mathbb{C}^3$ as a supplied parameter,
        consistent with the selection rule $\Sigma_c(n_3) = 3$, and reports a threshold rank
        $r_{\mathrm{eff}}^{1\%} = 3$ for every pair tested at $q \in \{61, 101, 151\}$, stating
        that this is a finite-data value of its protocol and not an invariant; alongside it O28
        reports a participation-ratio value $r_{\mathrm{eff}}^{\mathrm{PR}} = 8/3$.
        Paper O29~\cite{Beau2026a33} supplies finite-precision diagnostics compatible with an
        adjoint carrier but identifies no object: it states that no calculation there identifies
        these objects, no subspace of $\Heff$ is identified with $\Vr$, and the measured ranks are
        ranks in $\mathrm{End}(\Heff)$, not in $\mathrm{End}(\Vr)$.
        The value $d_\rho = 2$ is the minimality selection of O26~\cite{Beau2026a30}, which also
        records that rank three cannot be inverted to $d_\rho = 2$.
        Paper O27~\cite{Beau2026a31} proves a real-linear factorisation through the admissible
        quotient $\mathcal{Q}$ (Theorem~4.1(i)), without equivariance.
        Its identification $H_{\mathrm{eff}}^{\text{O27}} \simeq \su(2)$, where
        $H_{\mathrm{eff}}^{\text{O27}} = \mathrm{Im}\,\Pi_{\mathrm{adm}} \cap \mathrm{Ntrl}$,
        holds only under O23's supplied carrier together with O27's hypothesis of admissible
        saturation (Theorem~4.1(ii)), which O27 states is assumed and derived nowhere; the same
        hypothesis is what makes that space three-dimensional over $\mathbb{R}$ in O27.
        It is a statement about O27's own object and not about the measured $\mathbb{C}^3$
        (Section~\ref{sec:conv}).

  \subsection{The Problem}

    Both sub-programmes produce a three-dimensional effective structure.
    Their common substrate is the same group $\Heis(\mathbb{Z}/q\mathbb{Z})$, but they access it
    through different aspects: the Q-series through the BFS ball geometry, the O-series through
    the representation theory of the Weil blocks.
    The question studied in this paper is:

    \begin{quote}
      \emph{Are the three admissible directions of $\ImH$ (O23) and the three spatial directions
      of the effective Lorentzian geometry (Q5b) the same three-dimensional structure, seen from
      two different vantage points?}
    \end{quote}

    \noindent
    We show below that this question has a precise formulation, admits falsifiability criteria,
    but is \emph{not yet settled} within the corpus.
    The present work does not prove the identification.
    It states the requirements an identification would have to meet: under the identification
    hypothesis of Section~\ref{sec:sym2}, and with the spatial sector carrying the adjoint action
    [ACT], it is unique up to one scalar by irreducibility (Lemma~\ref{lem:schur}); and it must
    satisfy a specific quadratic-symbol constraint
    (Criterion~\ref{crit:symbol}) on an effective operator that exists only under the hypotheses
    of~\cite{Beau2026q5b}.
    Verdict, in the typing of the programme's bridge discipline: the required identification is
    \emph{missing}, not ruled out.
    The only obstruction located in the corpus is conditional and concerns a different route:
    under the depth and weight hypotheses of Q5a~\cite{Beau2026q5a}, no common scalar
    normalisation of the published admissibility form produces a toric differential limit
    (Q5a Theorem~6.1); nothing there or elsewhere proves that no bridge can exist.

  \subsection{Conventions}
    \label{sec:conv}

    Throughout, $q$ is an odd prime and $G_q = \Heis(\mathbb{Z}/q\mathbb{Z})$.
    We write $\pi_c : V_q \to \Heff = \mathbb{C}^3$ for the admissible projection of the
    character-$c$ Weil block, following the notation of O25--O29.
    The Heisenberg Lie algebra is $\heis = \mathrm{span}_\mathbb{R}\{X, Y, Z\}$ with $[X,Y] = Z$
    and $[X,Z] = [Y,Z] = 0$.
    The Lie algebra $\su(2)$ is spanned by $\{e_1, e_2, e_3\}$ with
    $[e_i, e_j] = \varepsilon_{ijk} e_k$ and is identified with $\ImH$ via
    $e_1 \leftrightarrow \mathbf{i}$, $e_2 \leftrightarrow \mathbf{j}$, $e_3 \leftrightarrow \mathbf{k}$.

    One notational warning, since the corpus uses one symbol for two objects.
    In O27~\cite{Beau2026a31} the symbol $H_{\mathrm{eff}}$ denotes
    $\mathrm{Im}\,\Pi_{\mathrm{adm}} \cap \mathrm{Ntrl}$, which that paper calls real
    three-dimensional \emph{under its hypothesis of admissible saturation}, and not otherwise;
    in O28--O29~\cite{Beau2026a32,Beau2026a33} and in the present paper $\Heff$ denotes the
    \emph{complex} three-dimensional target $\mathbb{C}^3$ of the admissible projection $\pi_c$.
    The warning O27 issues is about \emph{its own} pair of spaces, not about ours: it states that
    no dimension count of that paper transfers between its $H_{\mathrm{eff}}$ and the
    complexification $H_{\mathrm{eff}} \otimes_\mathbb{R} \mathbb{C}$ of that same space, in
    either direction.
    A fortiori no count transfers to or from the measured $\mathbb{C}^3$ of O28, which O27 does not
    discuss; that these two three-dimensional spaces are related at all is part of what
    Hypothesis~[ID] would have to supply.
    O23~\cite{Beau2026a27} separately warns against writing $\su(2) \cong \Sym^2(\Vr)$ without
    distinguishing a real from a complex dimension.
    We keep the complex reading throughout, write $W_{\mathrm{sp}}$ for the real spatial sector and
    $W_{\mathrm{sp}} \otimes_\mathbb{R} \mathbb{C}$ for its complexification, and use
    $H_{\mathrm{eff}}$ with the source spelled out whenever the O27 object is meant.

  \section{Two Independent Notions of Dimension Three}

  \subsection{The Geometric Three: Carnot Grading and Spatial Directions}
    \label{sec:geom3}

    The Heisenberg group $\Heis(\mathbb{R})$ is a Carnot group of step~$2$.
    Its Lie algebra $\heis$ carries a graded decomposition $\heis = V_1 \oplus V_2$ where
    $V_1 = \mathrm{span}\{X, Y\}$ is the horizontal distribution of Carnot degree~$1$ and
    $V_2 = \mathrm{span}\{Z\}$ is the central direction of Carnot degree~$2$.
    The sub-Riemannian Laplacian is $\Delta_H = X^2 + Y^2$ (acting on $\Heis(\mathbb{R})$
    by left-invariant vector fields); it involves \emph{only} the degree-$1$ generators.
    The homogeneous dimension is $\Dhom = 2 \cdot 1 + 1 \cdot 2 = 4$, giving BFS ball growth
    $|B_n| \sim C n^4$ (Theorem~2.1 of~\cite{Beau2026q5b}).

    The effective co-metric extracted in Theorem~5.2 of~\cite{Beau2026q5b}, conditionally on its
    hypotheses [H-L] and [H-lift], is
    $g^{\mu\nu} = \mathrm{diag}(-A_\tau,\, A_H,\, A_H,\, 0_{z\text{-sector}})$ in the basis
    $(\tau, \tilde X, \tilde Y, \tilde Z)$, where:
    \begin{itemize}
      \item $A_\tau > 0$ governs the temporal sector,
      \item $A_H > 0$ governs the two horizontal directions (entering $\Delta_H$ at
      principal-symbol level); its positivity is stipulated inside [H-L], not derived,
      \item the central slot is \emph{degenerate at this order}: Theorem~5.2
      of~\cite{Beau2026q5b} puts $0$ there and defers a full-rank spatial block to sub-principal
      corrections (open problem Q5b-O2). A value $A_z = 2$ is asserted in Q8~\cite{BeauQ8} and
      relayed by~\cite{Beau2026q5b} outside that theorem; this paper does not use it, for the
      reasons given in Section~\ref{ssec:withdrawn}.
    \end{itemize}
    Theorem~6.1 of~\cite{Beau2026q5b} is a separate statement, on the Lorentzian signature, and is
    conditional on Theorem~5.2 together with a hyperbolicity hypothesis.
    Throughout this paper, $A_z$ denotes the coefficient of $k_Z^2$ that a full-rank spatial block
    would carry. It is treated here as an unevaluated parameter: no value for it is derived in this
    paper, and the value asserted downstream is not used (Section~\ref{ssec:withdrawn}).

    \begin{remark}[Horizontal vs.\ derived: two sub-strata of the spatial sector]
      \label{rem:substrata}
      The three spatial directions of $g^{\mu\nu}$ do \emph{not} form a homogeneous set from the
      sub-Riemannian viewpoint.
      The directions $\tilde X, \tilde Y$ are horizontal (Carnot degree~$1$) and enter the effective
      operator at the level of the principal symbol.
      The direction $\tilde Z = [\tilde X, \tilde Y]$ is derived (Carnot degree~$2$) and is invisible
      to the principal symbol of $\Delta_H$; it contributes to the effective metric only through
      sub-principal corrections.
      This two-level structure --- degree-$1$ directions contributing at leading order, degree-$2$
      direction contributing at sub-leading order --- will play a key role in the bridge construction
      of Section~\ref{sec:bridge}, where it corresponds precisely to the distinction between
      spin-weight $\pm 1$ and spin-weight $0$ in the triplet $\Sym^2(\Vr)$.
    \end{remark}

    The three-dimensional spatial sector of $g^{\mu\nu}$ is the direct sum
    $V_1 \oplus V_2 = \heis / \mathbb{R}\tau$, where $\tau$ denotes the projective temporal direction.
    The ``geometric~$3$'' is $\dim_\mathbb{R}(V_1 \oplus V_2) = 3$.

  \subsection{The Admissible Three: The Neutral Sector of a Supplied Spinor Carrier}
    \label{sec:adm3}

    The admissible ``3'' arises at the representation-theoretic level.
    Paper O23~\cite{Beau2026a27} proves that if the admissible neutral sector is carried by an
    irreducible $\mathrm{SU}(2)$-valued two-dimensional representation
    $V_\rho \cong \mathbb{C}^2$ (a supplied input), then its traceless anti-Hermitian sector is $\su(2) \simeq \ImH$, carrying exactly
    three independent
    directions (Theorem~3.1); O23's Proposition~3.4, with its Remark~3.5, notes that a count of
    generator axes can stop
    at two, so ``directions'' is the correct reading.
    The carrier selection and the identification of $\Sigma_c$ with $\dim \su(V_\rho)$ are open
    bridges; the saturation threshold $\Sigma_c(n_3) = 3$ is a supplied selection rule.

    At the fibre level, the admissible projection $\pi_c : V_q \to \Heff$ lands in
    $\Heff = \mathbb{C}^3$, consistently with $\Sigma_c(n_3) = 3$~\cite{Beau2026a32}.
    The per-pair covariance of O28, a $9 \times 9$ operator on $\mathrm{End}(\Heff)$ which that
    paper writes $\mathcal{C}_c$, has threshold rank $r_{\mathrm{eff}}^{1\%} = 3$, its three
    resolved modes filling $\Heff$ at that numerical threshold~\cite{Beau2026a32}; O29~\cite{Beau2026a33} reports that
    this rank is modal rather
    than invariant across pairs and primes, and that it is a rank in $\mathrm{End}(\Heff)$, not in
    $\mathrm{End}(\Vr)$.
    No source identifies $\Heff$ with $\Sym^2_\mathbb{C}(\Vr)$: O29 states that no calculation
    there identifies these objects and that no subspace of $\Heff$ is identified with $\Vr$, and
    the value $d_\rho = 2$ is O26's minimality selection~\cite{Beau2026a30}.
    Accordingly, that identification enters this paper as a hypothesis
    (Hypothesis~[ID] of Section~\ref{sec:sym2}), on the data supplied by O23, O26 and O27.
    The ``admissible~3'' is the complex dimension of $\Heff$, a parameter of the pipeline, with
    which the modal rank of the measured covariance numerically coincides; the two are distinct
    quantities.

    \begin{remark}[Three distinct ``3''s to be kept separate]
      \label{rem:three3s}
      The present paper distinguishes three a priori distinct occurrences of~$3$ in the programme:
      (i) $\dim_\mathbb{R}(\heis / \mathbb{R}\tau) = 3$ (geometric, from Carnot structure),
      (ii) $\dim_\mathbb{R}(\ImH) = 3$ (algebraic, from the neutral sector of the supplied
      spinor carrier, O23 Theorem~3.1, conditional on that carrier),
      (iii) $\dim_\mathbb{C}(\Heff) = 3$ with modal covariance rank $r_{\mathrm{eff}}^{1\%} = 3$
      (measured, from O28, audited in O29).
      These coincide numerically but need not coincide structurally, and (iii) is a measured
      dimension carrying no representation-theoretic content by itself.
      Paper O23~\cite{Beau2026a27} states the analogous separation (carrier selection versus
      observable identification) as open bridges within the O-series itself.
    \end{remark}

  \section{\texorpdfstring{The Algebraic Obstruction: $\su(2) \not\simeq \heis$}%
    {The Algebraic Obstruction: su(2) is not isomorphic to heis(3)}}
  \label{sec:obstruction}

  \begin{proposition}[Non-isomorphism]
    \label{prop:noniso}
    The Lie algebras $\su(2) \simeq \ImH$ and $\heis = \mathfrak{heis}_3$ are non-isomorphic.
  \end{proposition}

  \begin{proof}
    The algebra $\su(2)$ is semisimple (its Killing form is negative-definite); $\heis$ is
    nilpotent of class~$2$ (its derived ideal $[\heis, \heis] = \mathrm{span}\{Z\}$ is central and
    non-trivial).
    A semisimple Lie algebra cannot be nilpotent.
  \end{proof}

  \noindent
  Proposition~\ref{prop:noniso} shows that the admissible ``3'' (which lives in $\su(2)$-space)
  and the geometric ``3'' (which lives in $\heis$-space) cannot be identified by a Lie algebra
  isomorphism.
  Any bridge between them must therefore operate at a level where the two structures are
  \emph{compatible without being isomorphic}: this is precisely the level of representations.

  \section{Structural Constraints on Any Bridge}
  \label{sec:bridge}

  \subsection{The Symmetric Square Identification}
    \label{sec:sym2}

    On the carrier supplied by O23~\cite{Beau2026a27}, with the minimality selection
    $d_\rho = 2$ of O26~\cite{Beau2026a30}, we have $\Vr \simeq \mathbb{C}^2$ as an
    $\su(2)$-module: the unique irreducible spin-$\tfrac{1}{2}$ representation.
    O27's saturation hypothesis is not needed for this step: it concerns
    $H_{\mathrm{eff}}^{\text{O27}}$, the space from which, by O27's own statement, no dimension
    count transfers here (Section~\ref{sec:conv}).
    Its symmetric square $\Sym^2_\mathbb{C}(\Vr)$ then carries a canonical basis
    $\{e_+, e_0, e_-\}$ of weight eigenstates under the Cartan generator $J_3$:
    $J_3 e_{\pm} = \pm e_{\pm}$, $J_3 e_0 = 0$.

    The measured space $\Heff$ and the module $\Sym^2_\mathbb{C}(\Vr)$ are, at this point, two
    distinct objects: the first is the three-dimensional target of the admissible projection
    $\pi_c$, the second is built from the supplied carrier.
    Their identification is an additional claim, and no source supplies it
    (Section~\ref{sec:adm3}).
    We therefore isolate it:

    \begin{quote}
      \textbf{Hypothesis [ID].} There is an $\su(2)$-module structure on $\Heff$ and an
      isomorphism of $\su(2)$-modules $\Heff \simeq \Sym^2_\mathbb{C}(\Vr)$.
    \end{quote}

    \noindent
    The following proposition states what representation theory contributes once [ID] is granted,
    and nothing more.

    \begin{proposition}[Symmetric square identification: conditional form]
      \label{prop:sym2}
      Up to isomorphism, $\Sym^2_\mathbb{C}(\Vr)$ is the unique irreducible three-dimensional
      complex module of $\su(2)$, with highest weight~$1$, and it is isomorphic to the
      complexified adjoint module $\su(2) \otimes_\mathbb{R} \mathbb{C}$.
      Consequently, under Hypothesis~[ID], $\Heff$ is an irreducible $\su(2)$-module of highest
      weight~$1$, and its weight decomposition under any choice of Cartan generator is
      $\{+1, 0, -1\}$, each weight space of dimension one.
    \end{proposition}

    \begin{proof}
      For the first assertion: the irreducible complex modules of $\su(2)$ are classified by
      highest weight, with $\dim V_j = 2j + 1$, so dimension three forces $j = 1$ and determines
      the module up to isomorphism.
      With $\Vr \simeq \mathbb{C}^2$ the spin-$\tfrac{1}{2}$ module, Clebsch--Gordan gives
      $\Vr \otimes \Vr = \Sym^2(\Vr) \oplus \wedge^2(\Vr)$ with
      $\dim \Sym^2(\Vr) = 3$ and $\wedge^2(\Vr)$ the trivial module, so
      $\Sym^2(\Vr) \simeq V_1$; and the complexified adjoint module of $\su(2)$ is the spin-$1$
      module of $\mathfrak{sl}_2(\mathbb{C})$, hence isomorphic to it.
      The second assertion is the transport of these facts along the isomorphism granted by [ID];
      the weight decomposition of $V_1$ is $\{+1,0,-1\}$ with one-dimensional weight spaces.
    \end{proof}

    \paragraph{What the proposition does not do.}
      Proposition~\ref{prop:sym2} derives no identification.
      In particular it does not follow from the O28 covariance data: the per-pair outer products
      $M_j = w_j w_j^{\mathsf T}$ formed from $w_j = \pi_c(v_j) \in \Heff$ lie in the symmetric
      part of $\mathrm{End}(\Heff)$, of complex dimension six, whereas
      $\Sym^2_\mathbb{C}(\Vr)$ has complex dimension three; the two spaces must not be conflated.
      Version~1.3 of this paper inferred [ID] from O29 and reported $d_\rho = 2$ as an O29 result.
      Neither attribution is retained (Section~\ref{ssec:withdrawn}).

  \subsection{Grading Correspondence}
    \label{sec:grading}

    The spin-weight grading on $\Sym^2(\Vr) = \mathrm{span}\{e_+, e_0, e_-\}$ decomposes the three
    basis states according to their $J_3$-eigenvalue $m \in \{+1, 0, -1\}$.
    The Carnot grading on $\heis$ decomposes the generators according to their homogeneous degree
    $d \in \{1, 2\}$.
    These are distinct gradings on distinct spaces, but they are compatible in the following sense.

    \begin{remark}[Grading correspondence]
      \label{rem:grading}
      The pair $(m, d) = (\pm 1, 1)$ (spin-weight $m = \pm 1$, Carnot degree $1$) corresponds to the
      horizontal generators $\tilde X, \tilde Y$, and the pair $(m, d) = (0, 2)$
      (spin-weight $m = 0$ within the $j = 1$ module, Carnot degree $2$) corresponds to the central
      generator $\tilde Z = [\tilde X, \tilde Y]$.
      The pattern is: the two generators entering the principal symbol of $\Delta_H$ (at leading order)
      have $|m| = 1$, while the generator entering via sub-principal corrections (at sub-leading order)
      has $m = 0$.
      This grading correspondence is not a Lie algebra isomorphism (Proposition~\ref{prop:noniso}), but
      it shows that the two sub-strata of the spatial sector identified in Remark~\ref{rem:substrata}
      map to the two spin-weight sectors of the spin-$1$ triplet in a grading-compatible way.
    \end{remark}

  \subsection{Rigidity of Any Bridge: Schur's Lemma}
    \label{sec:schur}

    The irreducibility of $\Sym^2(\Vr)$ as an $\su(2)$-module has an immediate and strong consequence
    for the possible forms of any bridge to the spatial sector.

    \begin{lemma}[Rigidity via Schur]
      \label{lem:schur}
      Let $W$ be a three-dimensional real vector space equipped with an $\su(2)$-action that decomposes
      as $W \simeq \mathfrak{su}(2)$ (adjoint representation, i.e.\ spin-$1$ over $\mathbb{R}$).
      When $W$ is taken to be the real spatial sector, that assumption on the target is the
      hypothesis isolated as [ACT] in Remark~\ref{rem:status}, which no source supplies.
      Then $\mathrm{Hom}_{\su(2)}(\Sym^2(\Vr), W \otimes_\mathbb{R} \mathbb{C}) \simeq \mathbb{C}$.
      Consequently an $\su(2)$-equivariant $\mathbb{C}$-linear map
      $\varphi : \Sym^2(\Vr) \to W \otimes_\mathbb{R} \mathbb{C}$ is either the zero map or an
      isomorphism, and in the latter case it is determined up to one non-zero complex scalar.
      Nothing here excludes $\varphi = 0$, and nothing here makes the scalar real or positive.
    \end{lemma}

    \begin{proof}
      The module $\Sym^2(\Vr)$ is irreducible of dimension~$3$ over $\mathbb{C}$.
      The complexification $W \otimes_\mathbb{R} \mathbb{C}$ of the real spin-$1$ representation is
      also irreducible of dimension~$3$ over $\mathbb{C}$ (the complexification of the adjoint
      of $\su(2)$ is the spin-$1$ representation of $\mathfrak{sl}_2(\mathbb{C})$).
      By Schur's lemma the space of equivariant maps between two isomorphic irreducible modules is
      one-dimensional, so every non-zero element of it is an isomorphism and any two differ by a
      non-zero scalar in $\mathbb{C}$.
    \end{proof}

    \begin{remark}[Consequence for the bridge]
      \label{rem:rigidity}
      Lemma~\ref{lem:schur} implies that, under Hypothesis~[ID], \emph{there is at most one
      equivariant bridge up to one scalar} between $\Heff$ and a three-dimensional adjoint module.
      That the real spatial sector \emph{is} such a module is the separate hypothesis [ACT] of
      Remark~\ref{rem:status}, which no source supplies.
      Uniqueness is guaranteed by irreducibility; existence, that is non-vanishing of $\varphi$, is
      not, and is the open question.
      A \emph{positive real} normalisation is not a consequence of any of this, and comparing forms
      does not deliver it either.
      Asking only that $\varphi$ carry a positive-definite Hermitian form to a positive-definite
      one constrains nothing, since $\varphi^* h(v,v) = |\lambda|^2 h(v,v)$ is positive definite for
      every $\lambda \neq 0$; asking for \emph{equality} with a specified form fixes the modulus
      $|\lambda|$ and leaves the phase free, a phase $e^{i\theta}$ acting trivially on that
      expression.
      Fixing the phase is a convention, to be stated wherever it is used; with the modulus fixed and
      the phase chosen, one may take the scalar to be $1$.
      Version~1.3 stated uniqueness ``up to positive scalar'' with neither the non-vanishing
      requirement nor this convention.
    \end{remark}

  \subsection{The Schr\"odinger Quadratic Form Lemma}
    \label{sec:qform}

    We now record the value of the $\su(2)$-Casimir on $\Sym^2(\Vr)$ and show that the Hermitian
    spatial form pulls back, under any non-zero equivariant bridge and if both coefficients are
    positive, to a positive-definite graded form on $\Sym^2(\Vr)$.

    The Casimir operator of $\su(2)$ on the spin-$j$ irreducible representation is
    $\mathcal{C} = j(j+1) \cdot \mathrm{Id}$.
    On $\Sym^2(\Vr)$ (spin $j = 1$), this gives $\mathcal{C}|_{\Sym^2} = 2 \cdot \mathrm{Id}$.

    \begin{lemma}[Schr\"{o}dinger quadratic form]
      \label{lem:qform}
      Write $W_\mathbb{C} := W_{\mathrm{sp}} \otimes_\mathbb{R} \mathbb{C}$ and let
      $\varphi : \Sym^2(\Vr) \to W_\mathbb{C}$
      be any $\su(2)$-equivariant isomorphism onto it: unique up to one scalar by
      Lemma~\ref{lem:schur}.
      The complexification is required, the source being a complex three-dimensional module and
      $W_{\mathrm{sp}}$ a real three-dimensional space.
      Then:
      \begin{enumerate}
        \item The quadratic form $Q_{\mathcal{C}}(v) := \langle v, \mathcal{C} v \rangle = 2\|v\|^2$
        on $\Sym^2(\Vr)$ is positive-definite of rank~$3$.
        \item If $A_H > 0$ and $A_z > 0$, then under any equivariant bridge $\varphi$ the pullback
        $\varphi^* h_{\mathrm{sp}}$ is positive-definite of rank~$3$, where $h_{\mathrm{sp}}$ is the
        Hermitian extension of the real spatial block
        $g_{\mathrm{sp}} = \mathrm{diag}(A_H, A_H, A_z)$ of $g^{\mu\nu}$ defined
        in~\eqref{eq:herm} below.
        The Hermitian extension is the one to use here: the $\mathbb{C}$-bilinear extension of a
        real form is never positive definite on a complex space, since it sends $iv$ to
        $-g(v,v)$.
        \item Fix the adapted isomorphism $\varphi_0$ that sends the weight vectors
        $(e_+, e_0, e_-)$ to the graded directions of $W_\mathbb{C}$ as in
        Remark~\ref{rem:grading}, normalised so that $\varphi_0^* h_{\mathrm{sp}}$ acts by $A_H$ on
        $\{e_+, e_-\}$ and by $A_z$ on $\{e_0\}$, and write $\varphi = \lambda\varphi_0$ with
        $\lambda \in \mathbb{C}^\times$, which is the freedom Lemma~\ref{lem:schur} leaves.
        Then
        \begin{equation}
          \varphi^* h_{\mathrm{sp}} = \lambda_H \cdot Q_{\pm} + \lambda_z \cdot Q_0,
          \qquad \lambda_H = |\lambda|^2 A_H, \quad \lambda_z = |\lambda|^2 A_z,
        \end{equation}
        where $Q_{\pm}$ is the projection onto the spin-weight $\pm 1$ eigenspace $\{e_+, e_-\}$ and
        $Q_0$ the projection onto the spin-weight $0$ eigenspace $\{e_0\}$; both coefficients are
        positive under the same proviso.
        The transport multiplies by $|\lambda|^2$, since
        $(\lambda\varphi_0)^* h(v,w) = h(\lambda\varphi_0 v, \lambda\varphi_0 w)
        = |\lambda|^2\, h(\varphi_0 v, \varphi_0 w)$; the reciprocal factor would belong to the
        transport in the other direction, along $\varphi^{-1}$.
        In particular the pullback depends on $\varphi$ only through $|\lambda|$, the phase acting
        trivially, which is why the normalisation discussion of Remark~\ref{rem:rigidity} concerns
        a modulus.
      \end{enumerate}
    \end{lemma}

    \begin{proof}
      Item~(1) is immediate from $\mathcal{C} = 2 \cdot \mathrm{Id}$ and $\dim \Sym^2(\Vr) = 3$.
      For items~(2) and~(3), fix first the real structure used throughout.
      The space $W_\mathbb{C} = W_{\mathrm{sp}} \otimes_\mathbb{R} \mathbb{C}$ carries the
      conjugation $\sigma(u \otimes \alpha) = u \otimes \bar\alpha$, an antilinear involution
      whose fixed-point set is exactly $W_{\mathrm{sp}}$; we write $\bar v := \sigma(v)$.
      Given the real symmetric form $g_{\mathrm{sp}}$ on $W_{\mathrm{sp}}$, its Hermitian extension
      is
      \begin{equation}
        \label{eq:herm}
        h_{\mathrm{sp}}(u \otimes \alpha,\, v \otimes \beta)
          := \bar\alpha\,\beta\; g_{\mathrm{sp}}(u, v),
        \qquad\text{equivalently}\qquad
        h_{\mathrm{sp}}(v, w) = g_{\mathrm{sp}}^{\mathbb{C}}(\bar v, w),
      \end{equation}
      with $g_{\mathrm{sp}}^{\mathbb{C}}$ the $\mathbb{C}$-bilinear extension.
      It is sesquilinear, satisfies $h_{\mathrm{sp}}(w,v) = \overline{h_{\mathrm{sp}}(v,w)}$,
      restricts to $g_{\mathrm{sp}}$ on $W_{\mathrm{sp}}$, and is positive definite exactly when
      $g_{\mathrm{sp}}$ is.
      Now $h_{\mathrm{sp}}$, equal to $\mathrm{diag}(A_H, A_H, A_z)$
      in the basis $(\tilde X, \tilde Y, \tilde Z)$, preserves the sub-strata of
      Remark~\ref{rem:substrata}:
      it acts by $A_H$ on $V_1 = \mathrm{span}\{\tilde X, \tilde Y\}$ (Carnot degree~$1$) and by $A_z$
      on $V_2 = \mathrm{span}\{\tilde Z\}$ (Carnot degree~$2$).
      The absence of cross terms between $V_1$ and $V_2$ in $g_{\mathrm{sp}}$ follows from
      invariance under the $\mathrm{U}(1)$ generated by $J_3$: weight spaces of distinct
      $J_3$-weight are orthogonal for any $J_3$-invariant Hermitian form, so such a form cannot mix
      the weight-$\pm 1$ sector with the weight-$0$ sector, and it is constant on the
      two-dimensional weight-$\pm 1$ sector.
      That is, $J_3$-invariance forces isotropy \emph{within the horizontal plane}; it does not by
      itself force $A_H = A_z$, since the weight-$0$ sector is a distinct isotypic component whose
      coefficient it leaves free.
      The $\mathrm{U}(1)$ argument is independent of the existence of the full bridge $\varphi$, but
      it requires the admissible regime to be $J_3$-invariant.
      That is a hypothesis of the present paper, isolated here as [INV], and it is not supplied by
      any source: \cite{Beau2026q5b} states its results in a ``homogeneous isotropic'' and
      ``homogeneous and quasi-isotropic'' regime, which is a restriction on the regime and not a
      named invariance hypothesis, and it contains no $J_3$ or $\mathrm{U}(1)$ statement at all.
      Under the grading correspondence of Remark~\ref{rem:grading}, $V_1$ maps to the spin-weight
      $\pm 1$ sector $\{e_+, e_-\}$ and $V_2$ maps to the spin-weight $0$ sector $\{e_0\}$.
      Under the grading correspondence the adapted $\varphi_0$ therefore pulls $h_{\mathrm{sp}}$
      back to $A_H$ on $\{e_+, e_-\}$ and $A_z$ on $\{e_0\}$, and
      $\varphi = \lambda\varphi_0$ multiplies both by $|\lambda|^2$, giving~(3).
      If $A_H, A_z > 0$, the pullback is positive-definite, giving~(2).
      Positivity is a hypothesis here, not an import: $A_H > 0$ is stipulated inside [H-L]
      of~\cite{Beau2026q5b}, and no positivity statement for $A_z$ is available from its
      Theorem~5.2, which puts $0$ in that slot.
    \end{proof}

    \begin{remark}[Status of Lemma~\ref{lem:qform}]
      \label{rem:status}
Item~(1) is structural: it is the statement that the Casimir acts on the spin-$1$ module as
$2 \cdot \mathrm{Id}$, and it holds under Hypothesis~[ID] alone.
Item~(3) and the full form of~(2) are conditional on five further premises: the existence of the
equivariant bridge $\varphi$ (Conjecture~\ref{conj:bridge} below); the existence of the effective
co-metric, which in~\cite{Beau2026q5b} is Theorem~5.2 under its hypotheses [H-L] and [H-lift],
with [H-L] not established there; the $\su(2)$-action on the real spatial sector that makes it the
adjoint module, which Lemma~\ref{lem:schur} requires of its target and which no source supplies
either, isolated here as [ACT]; the $J_3$-invariance [INV] of the admissible regime; and the
positivity of the two coefficients, of which $A_H > 0$ is stipulated inside [H-L] while $A_z$
receives no positivity statement from Theorem~5.2 of~\cite{Beau2026q5b}, which puts $0$ in that
slot; the value asserted downstream is not used here (Section~\ref{ssec:withdrawn}).
The lemma is therefore to be read as: \emph{if the bridge exists and those premises hold, it is
isometric in the graded sense} --- and in that case the spatial block is determined up to the two
constants $A_H$ and $A_z$, which the lemma does not evaluate.
    \end{remark}

  \subsection{The Bridge Conjecture}
    \label{sec:conj}

    The results above constrain the form of any equivariant bridge.
    We now state the remaining open question as a falsifiable conjecture.

    \begin{conjecture}[Bridge existence]
      \label{conj:bridge}
      Under Hypothesis~[ID], there exists a non-zero $\su(2)$-equivariant map
      \begin{equation}
        \varphi : \Sym^2(\Vr) \xrightarrow{\sim} W_{\mathrm{sp}} \otimes_\mathbb{R} \mathbb{C},
      \end{equation}
      where $W_{\mathrm{sp}} = \mathrm{span}_\mathbb{R}\{\tilde X, \tilde Y, \tilde Z\}$ is the
      three-dimensional real spatial sector of the effective co-metric, taken to carry the adjoint
      $\su(2)$-action [ACT].
      The map is required to be compatible with the Schr\"{o}dinger representation in the following
      sense.
      The compatibility is understood at the level of quadratic forms induced by the generators,
      not at the level of Lie algebra homomorphisms (which do not exist, by
      Proposition~\ref{prop:noniso}).
      Concretely: the quadratic form $Q_{\mathcal{C}} = 2\|\cdot\|^2$ on $\Sym^2(\Vr)$, when
      transported to $W_{\mathrm{sp}}$ via $\varphi$, matches the spatial block of
      $\sigma_2(L_{\mathrm{eff}})$, where $L_{\mathrm{eff}}$ is the effective operator whose
      existence is Theorem~5.2 of~\cite{Beau2026q5b} under its hypotheses [H-L] and [H-lift].
    \end{conjecture}

    \noindent
    The contract of the conjecture must be stated precisely, because both of its sides are
    conditional.
    Its source side is Hypothesis~[ID].
    Its target side exists only under [H-L], which \cite{Beau2026q5b} declares not established,
    the Lorentzian reading carrying the further hyperbolicity hypothesis of its Theorem~6.1; and
    that target carries the adjoint $\su(2)$-action [ACT] only by hypothesis, no source supplying
    it.
    Version~1.3 formulated the target as the spatial block of $\sigma_2(L_{\mathrm{eff}})$
    ``restricted to the admissible sector of $L^2(\mathbb{R})$ in the large-$q$ limit of Q5a''.
    That object is not supplied by Q5a and the formulation is withdrawn: in the published version
    of Q5a~\cite{Beau2026q5a}, under its depth hypothesis [H-depth] and its weight hypothesis
    [H-$w'$], every subsequential limit of the rescaled admissibility forms on toric targets with
    finitely many Fourier modes is the zero form, a bounded multiplication form containing no
    derivative term, or $+\infty$ (Q5a Theorem~6.1), the multiplication case having no
    differential symbol of positive order, hence no co-metric (Q5a Remark~6.2); in particular no
    such limit supplies a spatial second-order operator.
    That is a conditional obstruction to one route, under those two hypotheses and for common
    scalar normalisations of the published form on toric targets.
    It does not exclude a bridge obtained on a different filtration, with an anisotropic or
    non-scalar renormalisation, or on a non-toric target, and it proves nothing about the
    existence of $\varphi$.
    By Lemma~\ref{lem:schur}, if the map exists it is unique up to one scalar; by
    Lemma~\ref{lem:qform}, if it exists and the premises listed in Remark~\ref{rem:status} hold,
    the spatial block takes the graded form~(3).
    The conjecture is therefore an existence question about a missing identification, in the sense
    of verdict~3 of the programme's bridge discipline, and not a binary test whose negative answer
    would be a no-go.

  \section{Falsifiability Criterion}
  \label{sec:falsif}

  Lemma~\ref{lem:schur} removes the ambiguity about \emph{what} to test: under Hypothesis~[ID],
  and with the target carrying the adjoint action [ACT], a non-zero equivariant bridge is unique up
  to one scalar, so a single criterion suffices.
  It does not make the question binary in the sense of a no-go: a negative outcome bears on the
  tested realisation, as stated after Conjecture~\ref{conj:bridge}.

  \begin{criterion}[Symbol restriction to $\Sym^2(\Vr)$]
    \label{crit:symbol}
    Let $\sigma_2(L_{\mathrm{eff}})$ be the principal symbol of the effective operator
    $L_{\mathrm{eff}}$ on $\mathbb{R}_\tau \times \Heis(\mathbb{R})$, evaluated on the spatial block
    $(k_X, k_Y, k_Z)$.
    Via the grading correspondence of Remark~\ref{rem:grading}, identify the real spatial cotangent
    directions with a real basis of the image of $W_{\mathrm{sp}}$ inside
    $W_{\mathrm{sp}} \otimes_\mathbb{R} \mathbb{C} \simeq \Sym^2(\Vr)$: the horizontal pair
    $(k_X, k_Y)$ with $\tfrac{1}{\sqrt 2}(e_+ + e_-)$ and $\tfrac{i}{\sqrt 2}(e_+ - e_-)$, which
    span a real two-dimensional subspace of the weight-$\pm 1$ sector, and the central direction
    $k_Z$ with $e_0$.
    (The pairing $k_X \leftrightarrow e_+$, $k_Y \leftrightarrow e_-$ used in version~1.3 matched
    real cotangent coordinates with complex weight vectors and is corrected here.)
    Test whether:
    \begin{equation}
      \sigma_2(L_{\mathrm{eff}})\big|_{\Sym^2(\Vr)} = A_H (k_X^2 + k_Y^2) + A_z k_Z^2,
      \quad A_H, A_z > 0.
    \end{equation}
    The criterion is a test on the principal symbol of the continuum operator
    $L_{\mathrm{eff}}$ of~\cite{Beau2026q5b}.
    No such computation is performed in this paper or, to our knowledge, elsewhere in the corpus;
    Section~\ref{sec:numerical} computes a different object, namely compressions of the discrete
    Weil Laplacian.
    A failure of the identity would show that the tested realisation does not provide the bridge;
    it would not establish that no bridge exists.
  \end{criterion}

  \begin{remark}[What is computable, and what Section~\ref{sec:numerical} computes]
    \label{rem:compute}
    Section~\ref{sec:numerical} proves that the chirped discrete Fourier transform commutes with
    the discrete Weil Laplacian, $[F_c, L_{\mathrm{Weil}}] = 0$
    (Proposition~\ref{prop:commute}), and derives the consequence that a compression of
    $L_{\mathrm{Weil}}$ onto rows lying in distinct $F_c$-eigenspaces is block-diagonal
    (Corollary~\ref{cor:blockdiag}).
    Both are statements about the discrete operator on $\mathbb{C}^q$.
    They do not evaluate $\sigma_2(L_{\mathrm{eff}})$, and the transfer of a discrete compression
    to a continuum symbol coefficient is exactly the identification that Criterion~\ref{crit:symbol}
    presupposes and that no source supplies.
    Version~1.3 reported the isotropy condition as resolved, citing Q8~\cite{BeauQ8} for
    $A_z = C_{\mathfrak{su}(2)} = 2$ and Q10~\cite{Beau2026q10} with U1~\cite{Beau2026u1} for
    $A_H \to 2$.
    Those supports are no longer asserted here; see Section~\ref{ssec:withdrawn} for the reasons,
    which concern what this paper may rely on, not the editorial status of those papers.
  \end{remark}

  \section{What the O25 Checkpoints Measure}
  \label{sec:numerical}

  \subsection{Metaplectic symmetry of the discrete Weil Laplacian}
    \label{sec:metatplectic}

    The Weil Laplacian $L_{\mathrm{Weil}}$ on $\mathbb{C}^q$ for character $c$ is
    \begin{equation}
      L_{\mathrm{Weil}} = 4I - W_a - W_a^\dagger - W_b - W_b^\dagger,
    \end{equation}
    where $W_a$ is the cyclic shift $(W_a f)(k) = f(k-1 \bmod q)$ and
    $W_b = \mathrm{diag}(e^{2\pi i c k/q})_{k=0}^{q-1}$ is the character-$c$ phase operator.
    The chirped discrete Fourier transform $F_c$ is defined by
    \begin{equation}
      (F_c)_{jk} = \frac{1}{\sqrt{q}}\, e^{2\pi i c jk/q}, \qquad j,k \in \mathbb{Z}/q\mathbb{Z}.
    \end{equation}
    It is the metaplectic representative of the symplectic rotation $(x,\xi) \mapsto (-\xi, x)$
    in the character-$c$ Weil representation:
    \begin{equation}
      \label{eq:fc-conj}
      F_c W_a F_c^\dagger = W_b, \qquad F_c W_b F_c^\dagger = W_a^\dagger .
    \end{equation}
    (Version~1.3 printed the first relation as $F_c W_a F_c^\dagger = W_b^\dagger$; the sign is
    corrected here. The diagnostic of Section~\ref{ssec:diagnostic} checks both relations, the
    commutation $[F_c, L_{\mathrm{Weil}}] = 0$ and the deviation of the $\tilde F$ eigenvalue phases
    from multiples of $90^\circ$, and prints the residuals. The commutation is in any case
    unaffected by the sign, since $L_{\mathrm{Weil}}$ contains each generator together with its
    adjoint.)

    \begin{proposition}[Metaplectic symmetry]
      \label{prop:commute}
      For every odd prime $q$ and every $c$ with $\gcd(c,q) = 1$, the chirped DFT $F_c$ commutes with
      the Weil Laplacian:
      \[
        [F_c,\, L_{\mathrm{Weil}}] = 0.
      \]
    \end{proposition}

    \begin{proof}
      From the defining properties of $F_c$:
      \[
        F_c\, L_{\mathrm{Weil}}\, F_c^\dagger
        = 4I - F_c W_a F_c^\dagger - F_c W_a^\dagger F_c^\dagger
        - F_c W_b F_c^\dagger - F_c W_b^\dagger F_c^\dagger
        = 4I - W_b - W_b^\dagger - W_a^\dagger - W_a = L_{\mathrm{Weil}}.
      \]
      Hence $F_c L_{\mathrm{Weil}} = L_{\mathrm{Weil}} F_c$.
    \end{proof}

    \begin{corollary}[Block-diagonal structure]
      \label{cor:blockdiag}
      The matrix $L_{\mathrm{Weil}}$ is block-diagonal in every $F_c$-eigenbasis.
      Consequently, for any admissible basis $B_{\mathrm{eff}} \in \mathbb{C}^{3 \times q}$ whose rows
      lie in distinct $F_c$-eigenspaces, the restriction
      $\tilde{L} = B_{\mathrm{eff}}\, L_{\mathrm{Weil}}\, B_{\mathrm{eff}}^\dagger$ has no off-diagonal
      entries between sectors of different $F_c$-weight.
    \end{corollary}

    \noindent
    The mathematical content of Corollary~\ref{cor:blockdiag} is unaffected by this revision: it is
    the standard consequence that an operator commuting with $F_c$ preserves its eigenspaces.
    Its \emph{application} to the O25 data is withdrawn, on two counts.
    First, its hypothesis is not met by those data: the three stored rows are not
    $F_c$-eigenvectors, the off-diagonal entries of
    $\tilde F = B_{\mathrm{eff}} F_c B_{\mathrm{eff}}^\dagger$ being of order $q^{-1/2}$ rather
    than zero (a Fourier mode is mapped by $F_c$ to a vector concentrated at a single position,
    so the overlap of two such rows scales as $q^{-1/2}$).
    Second, even with the hypothesis met, the conclusion would concern off-diagonal entries of a
    compression of the discrete operator $L_{\mathrm{Weil}}$ on $\mathbb{C}^q$, not cross terms of
    $\sigma_2(L_{\mathrm{eff}})$: version~1.3 asserted the latter, which requires the missing
    identification of Criterion~\ref{crit:symbol}.
    What does hold for these data, exactly and for the reason given in
    Section~\ref{ssec:rayleigh}, is an arithmetic rule on the stored mode indices.

  \subsection{Compressions of the discrete Weil Laplacian, and their closed form}
    \label{sec:numval}
    \label{ssec:rayleigh}

    We compute compressions of $L_{\mathrm{Weil}}$ on the Q5a-O5 checkpoints of the O25
    pipeline~\cite{Beau2026a29}.
    For each prime $q \in \{61, 101, 151, 211\}$ and each tested character~$c$, the checkpoint
    stores the Gram--Schmidt admissible basis matrix $B_{\mathrm{eff}}$, a $q \times q$ complex
    matrix whose first three rows give the three admissible directions, together with the
    per-shell projections $\pi_c(v^{(k)})$ in $\mathbb{C}^3$, for $k$ in the fitting window.

    The computation proceeds in three stages.
    Stage~A computes the per-pair covariance $C_c = N^{-1}\sum_j w_j w_j^\dagger$
    in $H_{\mathrm{eff}} = \mathbb{C}^3$ and orders its eigenvectors.
    This $3 \times 3$ matrix is not the covariance O28 writes $\mathcal{C}_c$, which is a
    $9 \times 9$ operator on vectorised outer products in $\mathrm{End}(\Heff)$; the two carry the
    same letter in the corpus and are kept apart here.
    It is, on the other hand, the same object as the linear Hermitian covariance
    $\langle w w^\dagger \rangle$ whose spectrum O29~\cite{Beau2026a33} reports; the two differ only
    in aggregation, and the difference is worth stating because the figures look unrelated.
    The protocol here keeps the conjugate pairs separate, one covariance per pair.
    Per pair the normalised spectrum is $[1, 0.672, 0.666]$ at $(61,3)$, $[1, 0.920, 0.920]$ at
    $(61,5)$ and $[1, 0.500, 0.500]$ at $(61,7)$; these are three of the thirty conjugate pairs
    stored at $q = 61$, and the pairs differ widely.
    Taken over all thirty, the componentwise median of the per-pair spectra is $[1, 0.736, 0.732]$,
    which is the spectrum of no single pair, and over the
    five pairs that hold projections at $q = 29$ it is $[1, 0.512, 0.512]$; neither rounds to a
    figure O29 publishes.
    Pooling the projections instead, over the pairs that hold any, gives $[1, 0.7527, 0.7513]$ at
    $q = 61$ and $[1, 0.5584, 0.5584]$ at $q = 29$, which round to O29's $[1:0.8:0.8]$ and
    $[1:0.6:0.6]$.
    O29 states no aggregation in the remark that gives these two figures, but its protocol pools:
    elsewhere it reports a quantity as ``computed on the pooled Q5a-O5 checkpoints with the summary
    script'', and that script, \texttt{variety\_summary.py}, has a \texttt{pooled} function which
    concatenates the trajectories over the pairs of a prime.
    Forming $\langle w w^\dagger \rangle$ on its output gives the two spectra above.
    That last step is this paper's, not O29's: every covariance computed in O29's code is the
    $9 \times 9$ one on vectorised outer products, and no routine there forms this $3 \times 3$
    matrix.
    The agreement is at O29's one-decimal precision, and the pooled spectrum at $q = 61$ is not
    exactly degenerate, whereas O29's remark records $\lambda_2 = \lambda_3$.
    The per-pair figures are therefore not in disagreement with O29: they are a finer aggregation of
    the same measurement. Nothing in this paper depends on either.
    The obstacle to reproducing the pooled and median figures from the shipped bundle is coverage,
    not form: the pooled covariance is the mean of the per-pair covariances weighted by their
    projection counts, and the bundle carries the covariances, but it holds three of the thirty
    pairs at $q = 61$ and nothing at $q = 29$.
    Its projection counts sit beside it in \texttt{q7\_stage\_inputs\_meta.json}, which no script
    reads and which the digest does not cover.
    The same coverage limit puts the pair counts quoted here out of the bundle's reach.
    \texttt{--pooled --primes 29 61} in the diagnostic of Section~\ref{ssec:diagnostic} computes
    both from the checkpoints and prints with each the number of stored pairs, the number that hold
    projections and the number of projections pooled; the first two coincide at $q = 61$ and differ
    at $q = 29$, where $5$ of the $14$ stored pairs hold any.
    Stage~B checks that the $F_c$-eigenvalue phases of $\tilde{F} = B_{\mathrm{eff}} F_c B_{\mathrm{eff}}^\dagger$
    are multiples of $90^\circ$ (consistent with $F_c^4 = I$).
    Stage~C computes $\tilde{L} = B_{\mathrm{eff}} L_{\mathrm{Weil}} B_{\mathrm{eff}}^\dagger$,
    then conjugates it by the eigenbasis $U$ of the covariance $C_c$ of Stage~A, giving
    $L_{\mathrm{sw}} = U^\dagger \tilde L\, U$, and measures on $L_{\mathrm{sw}}$ the ratios
    $\rho_{Z|XY} = \max(|(L_{\mathrm{sw}})_{01}|, |(L_{\mathrm{sw}})_{02}|)/\bar\lambda$ and
    $|(L_{\mathrm{sw}})_{12}|/\bar\lambda$, where $\bar\lambda$ is the mean diagonal entry.
    The quantities historically named $A_z := (L_{\mathrm{sw}})_{00}$ and
    $A_H := \tfrac12\bigl((L_{\mathrm{sw}})_{11} + (L_{\mathrm{sw}})_{22}\bigr)$ are therefore
    diagonal entries of a compression of an operator on $\mathbb{C}^q$, read in the covariance
    eigenbasis.
    That second step matters: $U$ need not be a permutation, and whether it is free to rotate inside
    the plane of the two non-flat modes is decided by the spectrum of $C_c$ there, not by the
    entries of $\tilde L$ (Section~\ref{ssec:rotation}).

    \paragraph{The three stored rows are single Fourier modes.}
      Write $e_m(k) = q^{-1/2} e^{2\pi i m k/q}$ for the Fourier modes of $\mathbb{C}^q$.
      The diagnostic of Section~\ref{ssec:diagnostic} establishes that, for every tested pair, each
      of the three rows of $B_{\mathrm{eff}}$ is a single mode up to a phase, with spectral purity
      $1.000000$, the three indices being $\{0, -k, +k\}$ for an integer $k$ depending on
      $(q,c)$.
      This fixes every entry of $\tilde L$ in closed form.
      The generators act on modes by
      \begin{equation}
        \label{eq:mode-actions}
        W_a e_m = e^{-2\pi i m/q}\, e_m, \qquad W_b e_m = e_{m+c},
      \end{equation}
      the first because $e_m$ is an eigenvector of the cyclic shift, the second because
      multiplication by the character $e^{2\pi i c k/q}$ translates the Fourier index by $c$.
      Hence $\langle e_m, W_a e_m\rangle = e^{-2\pi i m/q}$ and, for $c \not\equiv 0$,
      $\langle e_m, W_b e_m\rangle = \langle e_m, e_{m+c}\rangle = 0$, so
      \begin{equation}
        \label{eq:rayleigh}
        \langle e_m, L_{\mathrm{Weil}}\, e_m\rangle
        \;=\; 4 - 2\cos\!\left(\frac{2\pi m}{q}\right).
      \end{equation}
      Two consequences follow, and they account for the whole of Table~\ref{tab:numerical}.
      At $m = 0$ the value is exactly $4 - 2 - 0 = 2$, for every $q$ and every $c$: it is the
      additive $4I$ term of $L_{\mathrm{Weil}}$ less the shift eigenvalue at the flat mode.
      At $m = \pm k$ the value is $4 - 2\cos(2\pi k/q) = 2 + 4\sin^2(\pi k/q)$.
      For the off-diagonal entries, $\langle e_m, L_{\mathrm{Weil}} e_n\rangle$
      $= -\langle e_m, W_b e_n\rangle - \langle e_m, W_b^\dagger e_n\rangle$ is non-zero exactly
      when $m - n \equiv \pm c \pmod q$, in which case its modulus is $1$.

    \begin{table}[ht]
      \centering
      \caption{What the compression $\tilde L = B_{\mathrm{eff}} L_{\mathrm{Weil}}
        B_{\mathrm{eff}}^\dagger$ contains on the O25 checkpoints.
        The three stored rows are the Fourier modes of indices $\{0, -k, +k\}$ (purity $1.000000$
        throughout). ``Flat'' is the entry at index $0$, ``mode $k$'' the common entry at indices
        $\pm k$, and ``closed form'' its value $4 - 2\cos(2\pi k/q)$ from~\eqref{eq:rayleigh};
        the largest discrepancy over all rows and pairs is below $10^{-14}$.
        The last column reports whether two stored indices differ by $\pm c \bmod q$, which is
        exactly when an off-diagonal entry of modulus $1$ appears; at $(61,3)$ this happens for two
        pairs of indices, at $(151,5)$ for one.
        Historically, $A_z$ and $A_H$ were read off this table after the Stage~A ordering; that
        ordering puts the flat mode first in eleven of the twelve pairs, and at $(q,c) = (101,3)$
        puts the mode $k = 5$ first, which is the whole origin of the published row
        $A_z = 2.096$, $A_H = \tfrac12(2.096 + 2.000) = 2.048$.}
      \label{tab:numerical}
      \smallskip
      \begin{tabular}{rrrccccc}
        \hline
        $q$ & $c$ & $k$ & flat & mode $k$ & closed form & $4\sin^2(\pi k/q)$ & cross term \\
        \hline
        61 &  3 &   3 & 2.000000 & 2.094729 & 2.094729 & 0.094729 & yes (two entries) \\
        61 &  5 &  29 & 2.000000 & 5.976176 & 5.976176 & 3.976176 & no  \\
        61 &  7 &   2 & 2.000000 & 2.042289 & 2.042289 & 0.042289 & no  \\
        101 &  3 &   5 & 2.000000 & 2.095974 & 2.095974 & 0.095974 & no  \\
        101 &  5 &  49 & 2.000000 & 5.991299 & 5.991299 & 3.991299 & no  \\
        101 &  7 &   3 & 2.000000 & 2.034730 & 2.034730 & 0.034730 & no  \\
        151 &  3 &   7 & 2.000000 & 2.084242 & 2.084242 & 0.084242 & no  \\
        151 &  5 &  73 & 2.000000 & 5.989188 & 5.989188 & 3.989188 & yes ($146 \equiv -c$) \\
        151 &  7 &   4 & 2.000000 & 2.027639 & 2.027639 & 0.027639 & no  \\
        211 &  3 &  10 & 2.000000 & 2.088020 & 2.088020 & 0.088020 & no  \\
        211 &  5 & 102 & 2.000000 & 5.989147 & 5.989147 & 3.989147 & no  \\
        211 &  7 &   5 & 2.000000 & 2.022128 & 2.022128 & 0.022128 & no  \\
        \hline
      \end{tabular}
    \end{table}

    \noindent
    Every entry of the compression $\tilde L$ \emph{in the stored basis} is thus accounted for
    by~\eqref{eq:rayleigh} and by the index-difference rule, with no free parameter and no appeal
    to the bridge. What the pipeline reports is $L_{\mathrm{sw}}$, which as a matrix coincides with
    $\tilde L$ at nine of the twelve pairs; at $(61,3)$ the two agree entry by entry in modulus but
    not in phase, and at $(101,3)$ and $(151,5)$ they differ outright, as
    Section~\ref{ssec:rotation} works out.
    In particular the two pairs at which version~1.3 reported anomalies, $(61,3)$ and $(151,5)$,
    are exactly the two pairs whose stored indices differ by $\pm c$: at $(61,3)$ the flat mode and
    the mode $-3$ differ by $c = 3$, giving $|\tilde L_{01}| = 1$ and hence the reported
    $\rho_{Z|XY} = 0.485$; at $(151,5)$ the modes $\pm 73$ differ by $146 \equiv -5$, which makes the
    restriction of $\tilde L$ to their plane non-scalar, and the $0.057$ printed in that row follows
    from that together with the covariance rotation, as Section~\ref{ssec:rotation} works out.
    Neither is a finite-size or finite-coverage effect.
    Since $A_H$ and the flat-mode value are two members of the one-parameter family
    \eqref{eq:rayleigh}, their difference is
    \begin{equation}
      \label{eq:gap}
      A_H - 2 \;=\; 2\bigl(1 - \cos(2\pi k/q)\bigr) \;=\; 4\sin^2(\pi k/q),
    \end{equation}
    an identity, not a measurement of a metric anisotropy.

    \begin{remark}[The flat-mode entry, and its coincidence with the $\mathfrak{su}(2)$ Casimir]
      \label{rem:az2}
      The flat-mode entry of the compression in the stored basis, $\tilde L_{00}$, equals $2$
      exactly, not approximately, at all twelve pairs.
      What depends on the Stage~A ordering is where that entry appears in the reported matrix
      $L_{\mathrm{sw}}$: at eleven pairs the flat mode is ordered first and
      $(L_{\mathrm{sw}})_{00} = 2$, while at $(q,c) = (101,3)$ the flat mode carries the
      \emph{smallest} covariance eigenvalue and is ordered last, so the $2$ appears in the last
      diagonal slot (Section~\ref{ssec:rotation}).
      By~\eqref{eq:rayleigh} this is the flat-mode Rayleigh quotient $4 - 2 - 0$ of
      $L_{\mathrm{Weil}} = 4I - W_a - W_a^\dagger - W_b - W_b^\dagger$: the $4$ of the additive
      normalisation, less $2$ from the shift eigenvalue at $m = 0$, less the vanishing character
      average.
      It is independent of $q$ and of $c$ by construction, and it is not a limit.
      The $\mathfrak{su}(2)$ Casimir does act on a spin-$1$ module as
      $j(j+1)\cdot\mathrm{Id} = 2\cdot\mathrm{Id}$, so the two numbers agree; but the agreement is
      not normalisation-invariant and therefore carries no evidential weight.
      Under $L_{\mathrm{Weil}} \mapsto \tfrac12 L_{\mathrm{Weil}}$ the entry becomes $1$ while
      $j(j+1) = 2$ is untouched, and every ratio of entries is unchanged; under
      $L_{\mathrm{Weil}} \mapsto L_{\mathrm{Weil}} - 2I$ it becomes $0$.
      Version~1.3 read this coincidence as a spectral signature of the bridge, and described the
      $(101,3)$ entry as a finite-window correction.
      Both readings are withdrawn (Section~\ref{ssec:withdrawn}).
      What survives is the identity~\eqref{eq:rayleigh} and the observation that any test of
      Criterion~\ref{crit:symbol} must compare normalisation-independent quantities.
    \end{remark}

    \begin{remark}[The cross term at $(q,c) = (61,3)$ is an index coincidence]
      \label{rem:finite}
      The cross-term ratio $\rho_{Z|XY} = 0.485$ at $(q,c)=(61,3)$ is not an outlier and not a
      finite-size artefact of the short window, as version~1.3 stated.
      The stored indices there are $\{0, -3, +3\}$ and the character is $c = 3$, so the flat mode
      differs from \emph{both} non-flat modes by $\pm c$ ($0 - (-3) = 3$ and $0 - 3 = -3$); by the
      index-difference rule of Section~\ref{ssec:rayleigh} both corresponding off-diagonal entries
      have modulus $1$, and dividing by the mean diagonal entry $\bar\lambda \approx 2.06$ returns
      the reported value.
      The same configuration occurs at $(151,5)$, where $\pm 73$ differ by $146 \equiv -5$.
      At every other tested pair no index difference equals $\pm c$ and the off-diagonal entries
      vanish identically, not approximately.
    \end{remark}

    \begin{remark}[The $c=5$ high-energy sector is outside the bridge test]
      \label{rem:c5}
      For all tested primes, the character $c=5$ gives a mode entry close to $6$, against $2$ at
      the flat mode.
      The mechanism is now explicit: the stored index at $c=5$ is $k = 29, 49, 73, 102$ for
      $q = 61, 101, 151, 211$, that is $k \approx q/2$, and~\eqref{eq:rayleigh} gives
      $4 - 2\cos(2\pi k/q) \approx 4 + 2 = 6$ at the top of the spectrum of
      $L_{\mathrm{Weil}}$.
      So the three stored rows for $c=5$ sit in the high-energy part of the spectrum, whereas for
      $c \in \{3,7\}$ the stored index satisfies $k/q \le 0.05$ and the entries sit near the
      bottom.
      Read in the direction the corpus needs, the observation is a statement about the pipeline's
      mode-selection rule, not about a physical sector: whichever sector one declares relevant,
      the selection of $k$ is what decides the value, and that selection is not derived anywhere.
      The $c = 5$ rows accordingly furnish no evidence either for or against
      Conjecture~\ref{conj:bridge}.
    \end{remark}

    \begin{remark}[Stage~B anomaly at $q=211$]
      \label{rem:q211}
      At $q=211$ (bfs\_frac$=0.11$), Stage~B does not pass for any of the three tested
      characters: the metaplectic overlap is $0.72$ ($c=3$), $0.60$ ($c=5$) and $0.38$ ($c=7$),
      all below the threshold $0.75$.
      Version~1.3 attributed this to low trajectory coverage, by analogy with the $(61,3)$
      anomaly. That attribution is not retained, and the following is what the quantity is.
      Stage~B does not test an entry of $\tilde F = B_{\mathrm{eff}} F_c B_{\mathrm{eff}}^\dagger$;
      it compares two unit vectors of $\mathbb{C}^3$, the dominant eigenvector of the covariance
      $C_c$ and the eigenvector of $\tilde F$ whose eigenvalue phase is closest to zero.
      Two features of $\tilde F$ hold on these data.
      The first follows from Section~\ref{ssec:rayleigh}: since $F_c$ maps a Fourier mode to a
      vector supported at a single position, every entry of $\tilde F$ has modulus exactly
      $q^{-1/2}$, so $\tilde F$ is $q^{-1/2}$ times a unimodular matrix.
      The second is measured rather than derived here: its eigenvalue phases agree with multiples of
      $90^\circ$ to within $10^{-8}$ degrees at all twelve pairs, consistent with $F_c^4 = I$; the
      diagnostic of Section~\ref{ssec:diagnostic} prints the deviation it measures per pair.
      The overlap itself, however, depends on the covariance direction and therefore on the
      trajectory: measured over the twelve pairs it ranges from $0.0000$ to $0.8610$, with no
      trend in $q$, and it does exceed the threshold at $(61,7)$ ($0.8610$) and $(151,3)$
      ($0.7864$) for bases of exactly this form.
      So the threshold is attainable, the failures are not explained by the basis structure, and
      this paper offers no explanation of the variation: what we can state is that no claim made
      here depends on Stage~B, whose outcome bears on the alignment of a measured covariance
      direction, not on the entries of Table~\ref{tab:numerical}.
    \end{remark}

    \begin{table}[ht]
      \centering
      \caption{The gap $\Delta(q) = 4\sin^2(\pi k(q)/q)$ between the mode entry and the flat-mode
        entry, for the stored indices $k(q)$ of the low-index characters.
        The values are identities, by~\eqref{eq:gap}, once $k(q)$ is read from the checkpoint; they
        are not measurements of a metric anisotropy.
        The $c = 3$ column is included because it is not monotone, which the version~1.3 reading
        did not report: $k(q) = 3, 5, 7, 10$ gives $k/q = 0.049, 0.050, 0.046, 0.047$.}
      \label{tab:asymp}
      \smallskip
      \begin{tabular}{rcccccc}
        \hline
        $q$ & $k$ ($c=7$) & $k/q$ & $\Delta$ ($c=7$) & $k$ ($c=3$) & $k/q$ & $\Delta$ ($c=3$) \\
        \hline
        61  &  2 & 0.0328 & 0.042289 &  3 & 0.0492 & 0.094729 \\
        101 &  3 & 0.0297 & 0.034730 &  5 & 0.0495 & 0.095974 \\
        151 &  4 & 0.0265 & 0.027639 &  7 & 0.0464 & 0.084242 \\
        211 &  5 & 0.0237 & 0.022128 & 10 & 0.0474 & 0.088020 \\
        \hline
      \end{tabular}

        {\small The extrapolated row at $q = 307$ printed in version~1.3 is removed: it had no
      asymptotic justification, and its value was placed in the $A_H$ column while being a value of
      $\Delta$.}
    \end{table}

    \noindent
    For $c = 7$ the stored index grows more slowly than $q$, so $k/q$ decreases and $\Delta$
    decreases with it; for $c = 3$ it does not, $k/q$ staying in $[0.046, 0.050]$ across the four
    primes.
    The four-point fit $\Delta \approx 0.364\,q^{-0.52}$ reported in version~1.3 is a fit to the
    sequence $4\sin^2(\pi k(q)/q)$ for the four selected indices $k = 2,3,4,5$: recomputing it from
    those indices alone returns the same exponent to two decimals.
    It therefore measures the growth of the selected index, and supports no statement about a
    spectral rate.

    \begin{conjecture}[Asymptotic isotropy]
      \label{conj:isotropy}
      Let $k(q,c)$ be the index of the non-flat mode selected by the pipeline at $(q,c)$.
      The gap between the two distinct entries of the compression vanishes,
      \[
        \Delta(q,c) = 4\sin^2\!\left(\frac{\pi k(q,c)}{q}\right) \longrightarrow 0
        \quad \text{as } q \to \infty,
      \]
      if and only if $\mathrm{dist}(k(q,c)/q, \mathbb{Z}) \to 0$.
      The conjecture to be settled is therefore a statement about the selection rule
      $q \mapsto k(q,c)$, and settling it requires deriving that rule, not fitting $\Delta$.
    \end{conjecture}

    \noindent
    The equivalence is immediate from~\eqref{eq:gap}, since $\sin^2(\pi x)$ vanishes exactly on the
    integers, and it separates the conjecture from the numerical support version~1.3 offered for
    it.
    A power-law fit to four values of $\Delta$ cannot establish the conjecture, because those four
    values are determined by four integers $k$ read off the same checkpoints; the fitted exponent
    is a property of that integer sequence.
    The $c = 3$ column of Table~\ref{tab:asymp} makes the point concretely: there $k/q$ does not
    decrease over the tested range, and neither does $\Delta$.
    What would make the conjecture testable is an independent determination of $k(q,c)$ from the
    breadth-first traversal and the Gram--Schmidt selection, which no paper of the corpus supplies
    and which the present paper does not attempt.

    \begin{corollary}[Status of the bridge identification]
      \label{cor:quasithm}
      The following is the status of the three claims this corollary carried in version~1.3.
      \begin{enumerate}
        \item The vanishing of off-diagonal entries of the compression $\tilde L$ holds, for the
        tested pairs, exactly when no two stored mode indices differ by $\pm c \bmod q$
        (Section~\ref{ssec:rayleigh}); two of the twelve pairs violate that condition.
        The corresponding statement about cross terms of
        $\sigma_2(L_{\mathrm{eff}})|_{H_{\mathrm{eff}}}$ is \emph{not} established: it requires the
        identification presupposed by Criterion~\ref{crit:symbol}.
        \item The flat-mode entry equals $2$ exactly and by construction
        (Remark~\ref{rem:az2}); its reading as a universal coefficient
        $A_z = C_{\mathfrak{su}(2)}$ is not retained.
        \item The claim $|A_H - A_z| < 0.05$ for $c \in \{3,7\}$, decreasing with $q$, is false as
        stated: at $c = 3$ the gap is $0.0947$ at $q = 61$ and $0.0842$ at $q = 151$, and it is not
        monotone in $q$.
        The exact value is $4\sin^2(\pi k/q)$ for the stored index $k$.
      \end{enumerate}
      Consequently no conclusion of the form $A_H = A_z = C_{\mathfrak{su}(2)} = 2$ is available
      from these data, and Conjecture~\ref{conj:bridge} is neither confirmed nor refuted by them.
      What remains is the conditional structure of Section~\ref{sec:bridge}: under
      Hypothesis~[ID], an equivariant bridge is unique up to one scalar, and if it exists the
      spatial block is graded as in Lemma~\ref{lem:qform}.
    \end{corollary}

    \begin{remark}[Three-case structure of the isotropy question]
      \label{rem:threecases}
      Criterion~\ref{crit:symbol} requires both the vanishing of cross terms and the equality
      $A_H = A_z$, for the symbol of the continuum operator.
      Neither condition is settled by the present data (Corollary~\ref{cor:quasithm}).
      For the equality, three genuinely distinct outcomes remain possible, with separate physical
      implications.
      \begin{enumerate}
        \item $A_H = A_z > 0$: the bridge exists and the effective spatial geometry is isotropic.
        The three spatial directions are metrically equivalent.
        This is the strongest conclusion: the two ``3''s are fully identified.
        \item $A_H \neq A_z$, both positive: the bridge exists but the effective spatial geometry is
        anisotropic.
        The three spatial directions are metrically distinct ($A_H$ for the horizontal, $A_z$ for the
        central).
        This is still a positive result: the diagonal structure is confirmed and the bridge is real,
        but the spatial metric is genuinely non-isotropic.
        \item $A_z \leq 0$ or structurally unstable: this realisation of the bridge fails, and the
        two ``3''s are not identified by it.
        Note that even this outcome would not prove that no bridge exists.
      \end{enumerate}
      Version~1.3 read Table~\ref{tab:numerical} as pointing to case~(1).
      That reading is withdrawn: the table contains compressions of the discrete operator, whose
      two distinct entries differ by the identity~\eqref{eq:gap}, so it discriminates between the
      three cases only through an identification it does not supply.
      No case is currently excluded.
      (Papers Q8~\cite{BeauQ8}, Q9~\cite{Beau2026q9} and Q10~\cite{Beau2026q10} cite this
      structure as ``Q7 Remark~6.8'': Q9 restates all three cases, Q8 and Q10 name Case~3 and
      Case~1 respectively. In the deposited version~1.3 and here
      it is Remark~\ref{rem:threecases}, the number 6.8 belonging to
      Corollary~\ref{cor:quasithm}. The numbering of
      this paper is unchanged in version~2.0.)
    \end{remark}

  \subsection{The covariance rotation, and the two pairs where it matters}
    \label{ssec:rotation}

    Two distinct things separate the reported matrix $L_{\mathrm{sw}}$ from the stored-basis
    compression $\tilde L$, and they must be kept apart because they pick out different pairs.

    \paragraph{The ordering.}
      The covariance eigenbasis $U$ is block-diagonal with respect to
      $\{\text{flat}\} \oplus \{e_{\pm k}\}$ at all twelve pairs, to numerical precision: the
      off-block entries of $C_c$ itself, relative to its largest eigenvalue, are below $10^{-14}$,
      and the diagnostic prints the figure it computes per pair.
      So the flat mode is not mixed into the others; what can change is the slot it occupies.
      At eleven pairs the flat mode carries the largest covariance eigenvalue and is ordered first.
      At $(q,c) = (101,3)$ it carries the \emph{smallest} (normalised eigenvalues
      $1$, $1$, $0.98758$) and is ordered last, which is the whole origin of the published row
      $A_z = 2.096$, $A_H = \tfrac12(2.096 + 2.000) = 2.048$.
      What is resolved there, by about one per cent, is the flat mode's separation from the plane;
      the plane itself is degenerate at $(101,3)$, as the next paragraph records.
      The two questions are distinct: this one is about the slot the flat mode occupies.

    \paragraph{The rotation inside the horizontal plane, and when it is harmless.}
      Whether $U$ is determined inside the plane of $e_{\pm k}$ is a property of the covariance
      spectrum, not of the Rayleigh quotients: equality of the two non-flat entries of $\tilde L$
      says nothing about $C_c$.
      The quantity to measure is the gap \emph{inside that plane}, not the gap between the second
      and third eigenvalues of the whole matrix: at $(101,3)$ the latter is the flat mode's
      separation and would mislead.
      Measured in the plane, with eigenvalues normalised by the largest of $C_c$, the gap is
      \emph{below $10^{-14}$} at ten pairs, $(101,3)$ included, one of them returning exactly zero,
      and it is resolved at two: $6.2 \times 10^{-3}$ at $(61,3)$ and $1.6 \times 10^{-2}$ at
      $(101,7)$.
      We quote a bound rather than the individual values, which sit at the level where the order of
      operations decides the last digits; the diagnostic of Section~\ref{ssec:diagnostic} prints
      them per pair, as it computes them.
      A gap of that size is a degeneracy \emph{to numerical precision}, which is all a computation
      on floating-point data can establish; we do not call it exact, no algebraic argument for it
      being offered here, and whether the coincidence is structural is a question about the
      pipeline that this paper does not settle.
      At those ten pairs the rotation of the plane is accordingly not determined by the data at
      the precision available: a different eigensolver may return a different basis of it.
      It is nevertheless without effect at nine of them, for a reason worth stating: there the
      restriction of $\tilde L$ to the plane is the \emph{scalar} operator
      $\bigl(4 - 2\cos(2\pi k/q)\bigr) I$, because no difference of the two indices equals
      $\pm c \bmod q$, and a unitary rotation of a plane on which an operator acts as a scalar
      returns the same scalar.
      No solver, ordering or seed can split those entries or create an off-diagonal term there.
      Exactly one pair has both an undetermined rotation and a non-scalar block, and it is
      $(q,c) = (151,5)$.
      Note that at the nine scalar-block pairs the conclusion does not depend on how the
      degeneracy is classified: scalar is scalar, so those entries are invariant whether the gap is
      exactly zero or merely below the noise floor.

    \paragraph{What happens at $(151,5)$.}
      There the two stored indices $\pm 73$ differ by $146 \equiv -5 \equiv -c$, so the restriction
      of $\tilde L$ to their plane is
      $\begin{pmatrix} 5.989188 & -1 \\ -1 & 5.989188 \end{pmatrix}$, which is not scalar, while
      the covariance plane is degenerate to well below $10^{-14}$, the diagnostic printing the
      figure it computes for the gap inside the plane at that pair.
      The rotation returned by the solver is not aligned with the stored modes, and it gives
      $L_{\mathrm{sw}} = \mathrm{diag}(2,\, 6.9527,\, 5.0257)$ with an off-diagonal entry of
      modulus $0.2677$ between the two non-flat slots.
      The value $5.989188$ survives only as the mean of the two split entries, by preservation of
      the trace, and the ratio $0.2677/\bar\lambda = 0.0574$ with $\bar\lambda = 4.6595$ is the
      number printed as $\rho_{XY} = 0.057$ in version~1.3.
      Those two entries and that off-diagonal term are therefore solver-dependent, and this is the
      only pair of which that is true.

    \paragraph{The column label.}
      Version~1.3's running text defined that column as the XY isotropy ratio
      $|\tilde L_{11} - \tilde L_{22}|/|\tilde L_{11}|$, on the matrix its own Stage~C had just
      introduced, namely the stored-basis compression
      $\tilde L = B_{\mathrm{eff}} L_{\mathrm{Weil}} B_{\mathrm{eff}}^\dagger$; its table caption
      instead called that same column a relative cross-term ratio.
      Neither description matches the number: on the stored-basis compression the formula gives $0$
      at every pair, both entries being equal there, while the quantity the script prints is the
      cross-term ratio $|(L_{\mathrm{sw}})_{12}|/\bar\lambda$ of the covariance-rotated matrix,
      $0.0574$ at $(151,5)$.
      Read as an isotropy ratio on that rotated matrix, the numerator is
      $|(L_{\mathrm{sw}})_{11} - (L_{\mathrm{sw}})_{22}| = 1.9270$; over the
      $|(L_{\mathrm{sw}})_{11}| = 6.9527$ the formula prescribes it gives $0.2772$, and over the
      mean $A_H = 5.9892$ the script's own isotropy variable uses it gives $0.3218$.
      Four readings, four distinct values, and the published column is the one its running text
      does not describe. The label is corrected here.

    \paragraph{The invariant content.}
      What does not depend on any of this is the spectrum of the compression on the
      three-dimensional stored span.
      It is $\{2,\, 4 - 2\cos(2\pi k/q),\, 4 - 2\cos(2\pi k/q)\}$ at the ten pairs where no
      difference of stored indices equals $\pm c \bmod q$, so there the closed
      form~\eqref{eq:rayleigh} gives the invariant spectrum outright.
      At the two pairs that do have such a difference it does not, and the spectrum is shifted by
      the off-diagonal entries: $\{0.632358,\, 2.094729,\, 3.462371\}$ at $(61,3)$, where the
      difference involves the flat mode, and $\{2,\, 4.989188,\, 6.989188\}$ at $(151,5)$, where it
      involves the two non-flat modes.
      The two criteria are independent: the index-difference rule says where the closed form
      departs from the invariant spectrum, and the covariance ordering says where the reported
      entries depart from the stored-basis entries. They coincide only at $(151,5)$.

    \paragraph{One published number is not accounted for.}
      The value $\rho_{XY} = 0.005$ printed at $(101,3)$ in version~1.3 is reproduced by no
      quantity the versioned script computes: at that pair the isotropy ratio is $0.0469$ with
      $A_H$ in the denominator and $0.0458$ with the caption's, and the cross-term ratios vanish to
      below $10^{-14}$.
      We record this rather than reconstruct it: the number's provenance is unknown, and no
      statement of this paper depends on it.

  \subsection{Reproducible diagnostic}
    \label{ssec:diagnostic}

    The statements of Sections~\ref{ssec:rayleigh} and~\ref{ssec:rotation} are checked by the
    script \texttt{code/q7\_mode\_diagnostic.py} of this repository.
    So that a reader with the repository alone can re-run it, the inputs it needs are shipped with
    it, in \texttt{code/data/q7\_stage\_inputs.npz} (37\,kB): per tested pair, the three stored
    basis rows, the $3 \times 3$ covariance $C_c$ of the per-shell projections, and the conjugate
    character pair, extracted from the O25 checkpoints
    \texttt{q\{61,101,151,211\}\_o25.npz} of the Q5a-O5 campaign.
    The bundle carries a SHA-256 digest in \texttt{code/data/SHA256SUMS}, which the script verifies
    before use and prints; \texttt{--from-checkpoints} recomputes everything from the full
    checkpoints instead, which are not redistributed here and which the loader materialises whole,
    so that path needs about 4.5\,GB of memory at $q = 151$ and 9.3\,GB at $q = 211$.
    The reader should note what this does and does not give: the bundle makes the statements of this
    section reproducible, since they depend only on those arrays, but it is a derived extract, so
    reproducing the extraction itself requires the campaign's checkpoints.
    The script prints, for each stored pair:
    the signed Fourier index of each of the three rows of $B_{\mathrm{eff}}$ and its spectral
    purity; the residuals of the two mode-action identities~\eqref{eq:mode-actions}; the measured
    diagonal entry of $\tilde L$, its closed form~\eqref{eq:rayleigh} and the difference; and, for
    each pair of rows, the measured off-diagonal modulus against the prediction of the
    index-difference rule.
    The Fourier convention is the one fixed in Section~\ref{ssec:rayleigh},
    $e_m(k) = q^{-1/2}e^{2\pi i m k/q}$, and the purity of a unit vector at index $m$ is
    $|(Fv)_m|^2/q$, equal to $1$ exactly when the vector is that single mode up to a phase.
    It also prints the normalised spectrum of $C_c$ with two gaps, the one between its second and
    third eigenvalues and the one inside the horizontal plane, keying its verdict to the second;
    the relative size of $C_c$'s off-block entries; the diagonal of the matrix after conjugation by
    the covariance eigenbasis together with the largest modulus of its off-diagonal entries; the
    spectrum of the compression; the residuals of the three metaplectic identities of
    Section~\ref{sec:metatplectic}, $F_c W_a F_c^\dagger = W_b$, $F_c W_b F_c^\dagger = W_a^\dagger$
    and $[F_c, L_{\mathrm{Weil}}] = 0$, together with the largest deviation of the $\tilde F$
    eigenvalue phases from a multiple of $90^\circ$; and a warning distinguishing the pairs where
    the plane is
    degenerate and the block is not scalar, so that the individual entries are solver-dependent,
    from those where the block is scalar and they are not (Section~\ref{ssec:rotation}).
    The \texttt{--pooled} flag instead reports, per prime, two aggregations: the covariance pooled
    over the conjugate pairs that hold projections, which is the one matching O29's published
    figures, and the componentwise median of the per-pair spectra, which matches neither.
    That path needs the checkpoints, for the reason given in Section~\ref{ssec:rayleigh}.
    On the four primes and three characters reported here the lowest purity is $1.000000$ and the
    largest discrepancy between measurement and closed form is below $10^{-14}$; every
    off-diagonal modulus of the compression in the stored basis matches the index-difference rule,
    being $1$ at $(61,3)$ (twice) and $(151,5)$ and $0$ elsewhere.
    The numerical agreement is reported here together with the derivation
    of~\eqref{eq:rayleigh}, so that the agreement is explained rather than merely observed.

    \clearpage
  \section{Discussion}
  \label{sec:disc}

  \subsection{What This Paper Establishes}

    The following results are proved in this paper, with their premises stated:

    \begin{itemize}
      \item The two occurrences of~$3$ in the programme arise via distinct mechanisms at distinct
      algebraic levels and cannot be identified by a Lie algebra isomorphism
      (Proposition~\ref{prop:noniso}). This is unconditional.
      \item Up to isomorphism, $\Sym^2(\Vr)$ is the unique irreducible three-dimensional complex
      module of $\su(2)$, isomorphic to the complexified adjoint module; under
      Hypothesis~[ID] the measured space $\Heff$ inherits that structure
      (Proposition~\ref{prop:sym2}). The identification itself is supplied, not derived.
      \item Under Hypothesis~[ID], an equivariant bridge from $\Heff$ to a three-dimensional
      adjoint module is either zero or an isomorphism, and in the latter case unique up to one
      non-zero complex scalar (Lemma~\ref{lem:schur}); reality and positivity of the scalar require
      the further requirement of Remark~\ref{rem:rigidity}.
      \item The $\su(2)$ Casimir acts on $\Sym^2(\Vr)$ as $2\cdot\mathrm{Id}$, and under a non-zero
      equivariant bridge, with both coefficients positive, the pullback of the Hermitian spatial
      form is graded as in Lemma~\ref{lem:qform}, with the premises listed in
      Remark~\ref{rem:status}. Invariance under the Cartan
      generator forces isotropy
      within the horizontal plane; it does not force $A_H = A_z$.
      \item The grading correspondence (Remark~\ref{rem:grading}) connects the Carnot degree-$(1,2)$
      decomposition of $\heis$ to the spin-weight-$(\pm 1, 0)$ decomposition of $\Sym^2(\Vr)$,
      which is why the $\tilde Z$ direction enters at sub-principal order. The two-dimensional
      horizontal distribution and the three-dimensional spatial sector are distinct objects and are
      kept distinct throughout.
      \item The chirped discrete Fourier transform commutes with the discrete Weil Laplacian for
      all $(q,c)$ with $\gcd(c,q) = 1$ (Proposition~\ref{prop:commute}), with the conjugation
      relations~\eqref{eq:fc-conj}.
      \item Every entry of the compression computed on the O25 checkpoints \emph{in the stored
      basis} is the Rayleigh quotient~\eqref{eq:rayleigh} of $L_{\mathrm{Weil}}$ on a single
      Fourier mode, with an off-diagonal entry exactly when two stored indices differ by
      $\pm c \bmod q$ (Section~\ref{ssec:rayleigh}, diagnostic in
      Section~\ref{ssec:diagnostic}). In particular the difference between the two distinct
      entries is the identity~\eqref{eq:gap}.
      \item The numbers the pipeline reports are those of the compression \emph{after} conjugation
      by the covariance eigenbasis, which is block-diagonal with respect to the flat mode and the
      horizontal plane at all twelve pairs. Two effects separate them from the stored-basis
      entries, and they are independent (Section~\ref{ssec:rotation}): the ordering of the flat
      mode, which differs only at $(101,3)$, and the rotation inside the horizontal plane, which is
      undetermined at the ten pairs whose covariance is degenerate to numerical precision inside
      that plane, but without effect wherever the restriction of the compression to the plane is
      scalar, that is at eleven of the twelve pairs. One pair, $(151,5)$, has both an undetermined
      rotation and a block that is not scalar, and it is the only pair whose reported entries are
      solver-dependent.
      \item The invariant content is the spectrum of the compression, which the closed
      form~\eqref{eq:rayleigh} gives outright at the ten pairs with no index difference equal to
      $\pm c \bmod q$, and which is shifted at the two that have one.
    \end{itemize}

    \noindent
    What the paper does \emph{not} establish is the identification of the two ``3''s.
    Its status is verdict~3 of the programme's bridge discipline: the required object, an
    equivariant bridge together with the module structure on the measured space, is not supplied by
    the audited corpus, and no result proves that it cannot exist.

  \subsection{Relation to Open Problem Q5b-O2}

    Open problem Q5b-O2 of~\cite{Beau2026q5b} asks for the extension of the effective co-metric
    to a full-rank spatial block including the $\tilde Z$-direction, via sub-principal corrections
    to the symbol of $L_{\mathrm{eff}}$.
    The present paper contributes a conditional statement of form, not a value: under
    Hypothesis~[ID] and the premises of Remark~\ref{rem:status}, if an equivariant bridge exists
    then the $\tilde Z$-direction contributes a term $A_z k_Z^2$ without cross terms with
    $(k_X, k_Y)$, and the horizontal pair enters isotropically.
    The coefficient $A_z$ is not evaluated here, and no computation of the sub-principal symbol of
    $L_{\mathrm{eff}}$ is performed in this paper.
    Version~1.3 reported the problem as resolved at the algebraic level, citing Q8~\cite{BeauQ8}
    for $A_z = C_{\mathfrak{su}(2)} = 2$, with Q10~\cite{Beau2026q10} and
    U1~\cite{Beau2026u1} for $A_H \to 2$, and treated the absence of cross terms as a theorem about
    the continuum symbol.
    None of these is asserted here; see Section~\ref{ssec:withdrawn}.
    Q5b-O2 is, from the point of view of this paper, open.

  \subsection{If the Identification Succeeds}

    If an equivariant bridge is exhibited and Criterion~\ref{crit:symbol} is met by the
    continuum symbol, the programme would hold a structural identification candidate between the
    three spatial directions of the effective geometry and the three-dimensional irreducible
    module of $\su(2)$.
    The three-dimensionality of space would then follow from the uniqueness of that module,
    conditionally on the supplied spinor carrier of O23~\cite{Beau2026a27} (Theorem~3.1, whose
    carrier selection is an open bridge), on O26's minimality selection~\cite{Beau2026a30}, and on
    the existence of the effective operator, which requires [H-L] of~\cite{Beau2026q5b}.
    That is a chain of three supplied premises, and it should be presented as such wherever the
    result is used.

  \subsection{If the Identification Fails}

    A failure of Criterion~\ref{crit:symbol} for a given realisation would show that this
    realisation does not identify $\Heff$ with the spatial sector: $\Heff$ would remain an internal
    representation space as far as that construction goes.
    It would bound the scope of the tested construction.
    It would not establish that the two ``3''s are structurally independent, which would require a
    proof that no admissible bridge exists; no such proof is available, and version~1.3's wording
    to the contrary is corrected.

  \subsection{Claims No Longer Asserted in Version 2.0}
    \label{ssec:withdrawn}

    This paper withdraws its own assertions and its own evidential uses.
    It takes no position on the editorial status of the papers cited below, whose claims are matters
    for those papers.
    Each item states what version~1.3 asserted, and why it is not asserted here.

    \paragraph{False statements about sources.}
      Version~1.3 stated that O29~\cite{Beau2026a33} identifies
      $\Heff \simeq \Sym^2(\Vr)$ as an $\su(2)$-module and establishes $d_\rho = 2$, and that the
      spin-$\tfrac12$ doublet is recovered as the Veronese square-root structure of the data.
      O29 states that no calculation there identifies these objects; $d_\rho = 2$ is O26's
      minimality selection~\cite{Beau2026a30}.
      Version~1.3 also stated that O27~\cite{Beau2026a31} proves unconditionally that any
      admissible morphism factors through $\su(2)$; O27 proves a real-linear factorisation through
      the admissible quotient, and the $\su(2)$ reading holds only under its saturation hypothesis.
      Version~1.3 attributed to Theorem~6.1 of~\cite{Beau2026q5b} a co-metric
      $\mathrm{diag}(-A_\tau, A_H, A_H, A_z)$ with $A_H, A_z > 0$ established there, and a
      convergence rate $O(q^{-1/2})$ to its Theorem~3.2.
      The co-metric is its Theorem~5.2 with $0$ in the central slot, $A_H > 0$ is stipulated inside
      [H-L], that theorem supplies no $A_z$, and its Theorem~3.2 is a Gromov--Hausdorff statement
      carrying no rate.
      (Q5b does relay $A_z = 2$ from Q8 elsewhere, in its abstract and in its Remark~5.4; the point
      here is only that Q7 attributed the value to the theorem it cited.)
      Version~1.3 located the target of Conjecture~\ref{conj:bridge} in a large-$q$ limit of
      Q5a~\cite{Beau2026q5a} that Q5a does not supply.

    \paragraph{Invalid arguments and invalid readings of computations.}
      The proof of Proposition~\ref{prop:sym2} in version~1.3 inferred the module structure from
      the O28 covariance, conflating the symmetric part of $\mathrm{End}(\Heff)$, of complex
      dimension six, with $\Sym^2(\Vr)$, of complex dimension three.
      Corollary~\ref{cor:blockdiag} was applied to data that do not satisfy its hypothesis, and its
      conclusion was transferred from a compression of $L_{\mathrm{Weil}}$ to cross terms of
      $\sigma_2(L_{\mathrm{eff}})$.
      The coincidence $A_z = C_{\mathfrak{su}(2)} = 2$ was read as a spectral signature of the
      bridge; it is the flat-mode Rayleigh quotient $4 - 2 - 0$ and is not normalisation-invariant
      (Remark~\ref{rem:az2}).
      The anomalies at $(61,3)$, $(151,5)$ and $q = 211$ were attributed to a short window, a
      high-energy sector and low coverage; they follow from the stored mode indices
      (Remarks~\ref{rem:finite}, \ref{rem:c5}, \ref{rem:q211}).
      The fit $|A_H - A_z| \approx 0.364\,q^{-0.52}$ was offered as evidence for asymptotic
      isotropy; it is a fit to $4\sin^2(\pi k(q)/q)$ on four selected indices
      (Conjecture~\ref{conj:isotropy}).
      Corollary~\ref{cor:quasithm} listed as proved or numerically supported three claims that are
      restated there with their actual status, one of them false as stated for $c = 3$.
      Lemma~\ref{lem:schur} was stated as giving uniqueness up to a \emph{positive} scalar without
      the requirement that makes the scalar real and positive, and without noting that the zero map
      is not excluded.

    \paragraph{Evidential uses of downstream results.}
      Version~1.3 reported the isotropy condition as resolved by Q8~\cite{BeauQ8}
      ($A_z = C_{\mathfrak{su}(2)} = 2$ by Casimir rigidity) and by Q10~\cite{Beau2026q10} with
      U1~\cite{Beau2026u1} ($A_H \to 2$ by spectral universality), and treated
      Criterion~\ref{crit:symbol} as thereby conditionally confirmed.
      This paper no longer uses those results as support.
      Two reasons are internal to the present revision and suffice for it.
      First, those statements concern the same two numbers that Section~\ref{ssec:rayleigh} shows to
      be members of the family~\eqref{eq:rayleigh}, so they cannot corroborate an identification
      that the family does not supply.
      Second, Q8's derivation fixes its scalar by matching $A_H$, and Q10's absolute bound is
      obtained by inserting $A_z = 2$ from Q8, so using them here would make this paper's support
      circular through its own Remark~\ref{rem:az2}.
      Q11~\cite{Beau2026q11} imports the same normalisation for the temporal coefficient; it is not
      used here either.

  \section{Conclusion}

  The programme exhibits two structural outputs equal to the integer~$3$: the Carnot spatial sector
  of $\Heis(\mathbb{R})$, whose effective co-metric is extracted in~\cite{Beau2026q5b} under its
  unestablished spatial limit hypothesis, and the admissible effective space
  $\Heff = \mathbb{C}^3$ measured in~\cite{Beau2026a32}, with the conditional spinor-carrier result
  of O23~\cite{Beau2026a27} supplying a three-dimensional neutral sector.
  These two ``3''s are numerically equal and algebraically distinct: $\su(2) \not\simeq \heis$.
  This paper states what an identification between them would require, and proves the requirements
  that follow from representation theory on the supplied carrier.
  Under Hypothesis~[ID], which identifies the measured space with the spin-$1$ module and which no
  source supplies, any equivariant bridge is either zero or an isomorphism unique up to one scalar,
  and, if that bridge is non-zero, the Hermitian spatial form pulls back to a graded form on the
  module, isotropic within the horizontal plane and determined by the two constants $A_H$ and
  $A_z$, neither of which is evaluated here.
  The open question is thereby given a precise contract rather than a resolution.

  The numerical section has been re-typed in version~2.0.
  The quantities it computes are compressions of the discrete Weil Laplacian onto three stored
  basis rows, and those rows are pure Fourier modes, so every entry is the Rayleigh quotient
  $4 - 2\cos(2\pi m/q)$: exactly $2$ at the flat mode, and $2 + 4\sin^2(\pi k/q)$ at the selected
  index~$k$, with an off-diagonal entry of modulus~$1$ exactly when two stored indices differ by
  $\pm c \bmod q$.
  The difference between the two distinct entries is therefore the identity
  $A_H - 2 = 4\sin^2(\pi k/q)$, and its vanishing in the large-$q$ limit is equivalent to
  $\mathrm{dist}(k/q, \mathbb{Z}) \to 0$, a property of the pipeline's mode-selection rule that no
  paper of the corpus derives.
  Accordingly the identification of the central value with the $\su(2)$ Casimir eigenvalue, the
  fitted convergence of the isotropy gap, and the transfer of the block-diagonality of a discrete
  compression to the cross terms of a continuum symbol are no longer asserted, and neither are the
  downstream supports version~1.3 invoked for them.

  What remains open is what was open before those supports were claimed: the bridge is a missing
  identification, in the sense of verdict~3 of the programme's bridge discipline.
  The corpus contains one conditional obstruction, on a different route and under the depth and
  weight hypotheses of Q5a~\cite{Beau2026q5a}: no common scalar normalisation of the published
  admissibility form yields a toric differential limit, so that route supplies no spatial
  second-order operator.
  It leaves open a bridge on a different filtration, under an anisotropic or non-scalar
  renormalisation, or on a non-toric target.
  Two concrete tasks would advance the question: a derivation of the mode-selection rule
  $q \mapsto k(q,c)$ from the traversal and the Gram--Schmidt step, which would make
  Conjecture~\ref{conj:isotropy} testable; and a computation of the sub-principal symbol of
  $L_{\mathrm{eff}}$ on the admissible sector, which is what Criterion~\ref{crit:symbol} actually
  tests and what no paper of the corpus has yet performed.

  \clearpage
  \bibliographystyle{plain}
  \bibliography{cosmochrony-bibliography}

\end{document}
